\begin{chapterbox}
    \chapter{Ein König auf Abwegen (2021)}
    \label{Ein König auf Abwegen (2021)}
    \az{Jahr 74}
    Die Rietgraskrone und die durch sie verliehene Verantwortung lasten schwer auf König Thorald. Sein Vertrauen in seinen manipulativen Ratgeber Ken Dorr schwindet. Seine geliebte Bäuerin hat er schon lange nicht mehr besucht. Zwei mächtige Schilde verstauben unter einer Klappe hinter dem Thron. Und hinter der Quelle des Likko versteckt sich ein langer Gang in ein Land voller Geächtete, Söldner, Sheriffs ... und einem fremden König.
\end{chapterbox}



\section{Thorald trifft eine Entscheidung}

\az{Jahr 74}

Die Nacht war bereits lange angebrochen und ihre finsteren Schatten schmiegten sich an die Mauern der Rietburg, als Thorald den Eingang erreichte. Er kniff die Augen zusammen, während er am Ewigen Feuer vorbeischlich, welches feuerrot in einer eisernen Schale vor dem großen Tor umhertollte. Er wollte nicht wahrhaben, dass kleine violette Flämmchen aufzischten, wenn er sich der Burg näherte. Er war doch keine Gefahr für das Königreich. Er war der König, bei Mutter Natur! Der König von Andor. Kein dahergelaufener Dieb.

„Man... man öffne die Tore! Ich bin heimgekehrt!“, entsprangen die krächzenden Worte den aufgerissenen Lippen des Thronerben. Die ganze Welt schwankte, als hätte man sie aus den Angeln gehoben. Thorald musste sich gegen das Tor stützen, um nicht dem feuchten Rietgras entgegenzustürzen. Schließlich gab er auf und sank daran herunter. Der Boden war matschig und kalt, aber es war so unendlich viel einfacher, sich der Erde zu ergeben, als sich gegen sie zu stemmen. Er hätte sich von der guten Gilda etwas Met für die Heimreise mitgeben lassen sollen, als Stärkung gegen die schändliche Schwerkraft. Aber so, wie er Gilda kannte, hätte sie ihm den Met glatt verweigert. Er lachte müde auf bei dem Gedanken. Sich dem König von Andor verweigern, eine einfache Tavernenwirtin!

„Wache! Ich erbitte um Einlass!“, lallte Thorald erneut, und endlich polterte etwas auf der anderen Seite der Burgmauer. Ein Gesicht blickt über die Brüstung. Armond, der Anführer der Rietgarde. Ausgerechnet heute musste Armond Nachtschicht haben, dieser pflichtbewusste Besserwisser.

„Mein König“, sprach Armond knapp, als er Thorald ins Innere der Burg führte. Er bot Thorald bereitwillig seinen Arm als Stütze an. Gleichzeitig bedachte er ihn mit einem Blick voller Abscheu und Mitleid, bei dem Thorald am liebsten auf der Stelle im Boden versunken wäre.

„Ich finde den Weg alleine“, blaffte Thorald den Gardisten an, „Ist ja nicht das erste Mal, dass ich... Ihr wisst schon.“

„Ja, das ist leider nicht das erste Mal“, gab Armond resigniert zurück. Er setzte bereits an, zu seinem Posten zurückzukehren, als er stutzte.

„Mein König, ist das Blut? Seid Ihr verletzt? Seid Ihr auf Kreaturen gestoßen? Ich sagte Euch doch, dass Ihr Euch von einer Leibgarde begleiten lassen sollt, wenn Ihr...“

„Schweigt stille, Ihr... Ihr La... La... Langeweiler. Ich wurde vom besten Schwertmeister ausgebildet, den dieses Königreich je ge... gesehen hat... da werde ich doch noch hoffentlich mit einigen ga... ga... garstigen Gors klarkommen können.“

Die Welt schwankte erneut und Thorald sackte auf ein Knie. Er wusste, dass eine Ausbildung bei Harthalt kein Garant dafür war, eine Begegnung mit Kreaturen unbeschadet zu überstehen – Harthalt selbst hatte sein Leben im Kampf gegen einen einzelnen Skral verloren, und war dabei noch nüchtern gewesen. Thorald wusste auch, dass es nicht geschickt war, sich ohne Begleitung zum Trinken zu schleichen. Aber in Begleitung von jemandem, der sich seines Königs schämte, wären seine Ausflüge noch viel unerträglicher, als sie es ohnehin schon waren. Und es war nicht so, als wäre das Land auf seine Anwesenheit im Thronsaal angewiesen, um am Laufen zu bleiben.

Dass er heute Nacht auf zwei grantige Gors gestoßen war, war bloß Pech gewesen. Pech, das ihn vielleicht sein Pferd gekostet hatte. Er sollte einen Boten aussenden, um das umliegende Rietland nach seiner schönen Stute abzusuchen. Mit etwas Glück war sie nur durchgegangen und unversehrt geblieben. Morgen, nahm er sich vor, morgen würde er das tun. Ach, wenn doch nur Thorn noch hier wäre. Der Bursche wusste, wie man mit Pferden umgeht. Und er hatte stets aufmunternde Worte für Thorald gehabt. Nun, zumindest so lange, bis Thorald ihn als nichts als einen dummen Bauernsohn mit einem Schwert bezeichnet hatte. Eine weitere Beziehung, die Thorald gekappt und es dann bereut hatte. Thorald schwelgte noch einige trunkene Augenblicke in Erinnerungen an Thorn, dann rappelte er sich ächzend auf.

„Mein König, wenn Ihr verletzt seid, so muss ich darauf bestehen, dass Ihr mich zum Heiler begleitet“, versuchte Armond erneut sein Glück.

„Lasst den alten Alduin doch schlafen“, befahl Thorald, „Mir geht es doch... so... so gut wie neugeboren.“

Armond schüttelte bloß den Kopf. Ein Kontrollblick verriet ihm allerdings, dass der König noch sämtliche Finger und Ohren besaß und dass die Flecken auf seinem Wams eher nach Wein als nach Blut aussahen. So ließ er ihn widerwillig ziehen.\bigskip







Thorald stapfte durch die Rietburg. Der Mond warf sein Licht auf die Abdrücke von Thoralds edlen Stiefeln im Matsch. Dreck spritzte an seine ohnehin schon besudelten Leinenhosen. Die Burg lag friedlich da. Die meisten Leute schliefen in Ruhe, während ihr König betrunken heimkehrte.

Eine Eule schrie in der Ferne, und Thorald zuckte zusammen. „Nur ein Eule, nur ein einfacher Flattervogel“, murmelte er zu sich selbst.

Langsames, rhythmisches Hämmern traf auf seine Trommelfelle. Hatte sich Warguth tatsächlich diese Tageszeit ausgewählt, um seine Schmiedekunst zu perfektionieren? Seitdem Wulfron die Schmiede an Warguth abgetreten hatten, war die Qualität der darin produzierten Schwerter laut Ken merklich gesunken – und Ken hatte ein Auge für so etwas. Das war aber noch lange kein Grund, bis mitten in die Nacht hinein verbissen weiterzuwerkeln. Thorald grummelte. Der Hall in seinen Ohren war alles andere als angenehm. Wie konnten die restlichen Bewohner der Rietburg auch nur ein Auge zutun dabei?

Tatsächlich, da war Warguth, tief in seine Arbeit versunken. Funken sprühten, während sein Schmiedehammer in regelmäßigen Intervallen auf den Amboss niederfuhr. Warguths schwarzer Hund lag neben ihm und hob den Kopf, als er Thorald vorbeistapfen sah. Ein tiefes Knurren entsprang seiner Kehle. Nicht einmal die Hunde mochten Thorald noch.

Nun fiel auch Warguth die Präsenz seines Königs auf und der junge Schmied hielt in seinem Werk inne. Er sah Thorald entgegen, wie dieser als elendes Häufchen den Trampelpfad entlangtorkelte. Kurz wirkte es, als wollte Warguth etwas sagen, doch dann wandte er sich bloß wieder seinem Schwert zu und hämmerte stirnrunzelnd weiter. Thorald empfand das Bedürfnis, etwas zu erwidern, doch die passenden Worte fielen ihm nicht ein und ganz abgesehen davon konnte er es kaum erwarten, in sein wohliges weiches Bett zu fallen und die Welt um ihn herum zu vergessen. Also zog er weiter.\bigskip







Zum Schlafen sollte Thorald so schnell nicht kommen, denn in seinen Gemächern erwartete ihn ein ganz und gar nicht gut gelaunter Ken Dorr. Ken hatte natürlich kein Auge zugetan und setzte sich mit verschränkten Armen auf, sobald Thorald durch die Tür stolperte.

„Thorald“, setzte Ken an, sein Tonfall von Missgunst und Unmut geradezu triefend. Thorald brachte eine wegwerfende Handbewegung zustande und brummelte etwas in seinen Bart hinein, während ihn noch mehr Schuldgefühle überkamen.

„Wir haben schon einige Male darüber geredet. Ich glaubte, mich klar ausgedrückt zu haben. Dein Verhalten ziemt sich nicht eines Königs. Alleine zum Umtrunk auszureiten und ohne Pferd und ohne Besinnung zurückstolpern zu kommen – was denkst du dir nur immer dabei?“

Thorald ignorierte Kens lauter werdende Stimme und schüttelte seine verdreckten Stiefel ab. Zudem wandte er seinem engsten Ratgeber den Rücken zu, damit dieser möglichst nicht die Blutflecken auf seinem Wams bemerkte.

„Währenddessen macht sich die ganze Rietburg Sorgen machen um dich, aber das hat dich noch nie gekümmert. Dir hätte wer-weiß-was alles zustoßen können! Was würde deine Mutter sagen? Und Brandur? Warum bist du nur so furchtbar unverbesserlich? Was muss ich tun, um zu dir durchzudringen?“

„Jetzt zünd‘ mal nicht gleich das ganze Rietgras an!“, fuhr Thorald zurück, „Allen anderen Andori ist es doch völlig egal, was mit mir geschieht. Meine Regen... Regent... Regentschaft zeichnet sich doch bloß dadurch aus, dass ich dumme Entscheidungen treffe, die sich schlecht auf das Königreich auswirken. Meinetwegen sind so viele Zwerge im Kampf um Cavern umgekommen. Die hassen uns jetzt noch mehr. Und es war ich, der um die Beachtung und Liebe der Bevölkerung buhlte und darum die Helden von Andor in den Norden sandte, um nicht in ihrem Schatten zu stehen. Wenn die nächste Trollhorde aus dem Gebirge hier aufkreuzt, sind wir ihr schutzlos ausgeliefert. Alle diese verlorenen andorischen Leben werden auf meine Kappe gehen! Selbst die Steu... Steuern waren keine gute Idee, die haben die Flussländler gegen uns aufgebracht. Ich wollte meinem Vater nacheifern und in die Legenden von Andor eingehen, aber inzwischen glaube ich... ich glaube, dass ich einfach nur ein Tor bin und nicht für dieses Amt geschaffen.“

Manchmal wünschte sich Thorald die Zeit zurück, in welcher er nur ein kleiner, naiver Junge gewesen war. Damals hatte er eine Selbstsicherheit besessen, die an Hochmut gegrenzt, ja, in manchen Momente diese Grenze gar überschritten hatte. Nichts hatte seine Selbstsicherheit damals trüben können. Natürlich hatte er nicht weniger Fehler gemacht als heute – er war es gewesen, dessen voreiliges Ausziehen der Rietgarde die Rietburg für ihre erste Eroberung schutzlos zurückgelassen hatte. Aber zumindest war er sich damals der Konsequenzen seiner Handlungen nicht bewusst gewesen. Und zumindest war er damals noch von fröhlichen Ratgebern und Lehrern umgeben gewesen. So vieles hatte sich verändert seit damals.

Thorald mühte sich vergeblich mit seinem verdreckten Wams ab, grunzte erschöpft und ließ dann endlich die Schwerkraft ihr Ding tun. So fiel er mehr neben Ken ins Bett, als dass er hineinstieg, und verteilte vergorenen Wein und vergossenes Gorblut darüber.

„Vielleicht sollte ich abdanken. Vielleicht kommt ja jemand und erhebt Anspruch auf den Thron. Irgendein entfernter Verwandter, der die Lage richten kann.“, murmelte er in seine Strohmatratze. Kurz fühlte er in seiner Brust den Stich der verfluchten Klinge des Schwarzen Ritters, mit der sein Onkel ihn einst hatte umbringen lassen wollen. Ob Hademar wohl Kinder hatte, die einen Anspruch auf den Thron Andors geltig machen und ihn aus diesem elenden Amt befreien konnten?

Thorald wusste nicht, ob Ken sein Gemurmel gehört hatte. Auf jeden Fall atmete Ken tief durch, seine Stirn glättete sich und er kratzte gedankenverloren an der Narbe an seiner rechten Wange herum. Von Gezeter und Standpauke keine Spur mehr.

„Mein König... du bist zu hart mit dir selbst“, sprach Ken fest und ließ sich neben Thorald nieder, „Unter unserer Führung ist das Königreich erblüht wie noch nie. Die Anstürme der Kreaturen haben endlich nachgelassen. Die Vorratskammern sind voll. Die diplomatischen Beziehungen mit Tulgor sind auf einem Höhepunkt. Das Volk liebt dich, Thorald. Du als König hast ihnen die Ära des Friedens gebracht, für den sie so lange kämpften. Du stärkst alle Personen in deiner Nähe. Du bist ein Held. Der Thron Andors ist sicher in deiner Hand.“

Thorald wälzte sich herum und blickte Ken an. Er wollte ihm widersprechen, wusste er doch sehr wohl, welchen Schaden seine voreiligen Entscheide immer wieder über das Reich gebracht hatten. Aber Kens Lobe fühlten sich stets so gut an. Vielleicht konnte er ja doch ein klein wenig Verantwortung für den momentanen Frieden übernehmen. Über acht Jahre davon. Das hatte es in der Geschichte Andors noch nie gegeben.

Aber nein, wenn er ehrlich mit sich war, dann waren die meisten guten Entscheidungen aus Thoralds Regentszeit Kens Kopf entsprungen. Ken war es, dem das Volk den Königsfrieden verdankte. Und Ken war noch nicht fertig mit ihm.

„Deine Probleme, Thorald, haben nicht mit deiner Regierung zu tun, sondern mit deiner Verantwortungslosigkeit. Der Druck der Rietgraskrone macht dir zu schaffen, natürlich, wem würde er nicht? Aber sich alleine aus der Burg zu schleichen wie ein ungezogener Bengel, das ist unter deiner Würde. Es muss doch bessere Möglichkeiten geben, mit deiner Frustration umzugehen. Wenn du diese nicht finden kannst, wirst du früher oder später daran zu Grunde gehen. Thorald, du musst stark bleiben. Wir stehen das durch.“

Thorald seufzte. Rekas Warnung schoss wieder einmal ihm durch den Kopf. Und so begann er, vor sich hin zu murmeln:

„Ken, du weißt, dass dein Rat mir von allen am teuersten ist. Ernsthaft. Wenn noch welche von Wulfrons Heldenbroschen übrig wären, würde ich dir höchstoffiziell eine verleihen für deine Dienste für das Königreich“

Kens Mund verzog sich zu einem breiten Lächeln. Thorald brauchte mehrere Anläufe, um die bösen Worte hervorzustoßen, die ihm auf dem Herzen lagen.

„Aber ich komme nicht um... umhin, zu befürchten, dass deine Worte leer sind. Dass du mich nicht um m... m... meiner selbst willen willst. Du hast doch bloß Angst davor, dass, wenn ich abdanke, ein anderer an meiner Stelle tritt. Einer, der von selbst klarkommt mit all den Schwierigkeiten des Königsamtes. Einer, der sich nicht so leicht kontrollieren lässt wie mich.“

Ken lächelte nun nicht mehr. Thorald ebensowenig.

„Ken... die alte Reka hat kürzlich mir eine Warnung überreicht. Dass mich einer manipuliere. Einer, der im Geheimen bereits die Rietgraskrone trage. Dass einer davon träume, an meiner Stelle auf dem Thron zu sitzen. Dies... dies kann nur auf dich zutreffen. Ich mag kein weiser König sein, aber ich weiß zumindest so viel: Du sollst... du sollst mich nicht kontrollieren. Und das kannst du nicht. Nicht mehr länger. Ich bin ein freier Mensch und dein Rat... dein Rat ist nur so viel... ein Rat halt. Lass mich dir zur Abwechslung einen Rat geben: Es geht dich nichts an, wenn ich mich in der Nacht einen Umtrunk gönne. Und es gehört dir nicht, mir zu widersprechen, wenn... wenn ich etwas... wenn ich die Krone zum Wohle des Reiches abtrete... ooo... Ken, mir ist kotzübel.“

Ken stand brüsk auf und starrte auf Thorald nieder, welcher immer noch in seinem verdreckten Wams dalag und vor sich hin murmelte. Kurz verharrte er regungslos, doch dann ließ er sich auf ein Knie fallen und flüsterte: „Mein König! Wie kannst du auch nur daran denken, dass es mich nach deiner Krone verlangte? Ich sehe doch, wie sehr du an diesem Los zu leiden hast. Diese Reka... wie alt ist sie nun schon? Einhundert Jahre? Mehr? Sie hat nicht mehr alle Tassen im Schrank, das ist es. Sieht hinter allem Intrigen und will vermutlich ihre eigene politische Macht mehren. Du kannst ihr nicht vertrauen. Mir schon. Deinem treuen Ken Dorr. Komm schon, Thorald, lass nicht zu, dass deine Zweifel dich übermannen.“

Thorald schüttelte bloß seinen Kopf und schluchzte erstickt etwas in sein Kissen hinein, von dem er nicht einmal selbst wusste, was es bedeuten sollte. Die Übelkeit in seinem Innern nahm stetig zu.

„Das reicht jetzt, Thorald“, sprach Ken mit einem bedeutungsschwangeren Unterton in der Stimme. „Du hast dich lange genug geplagt. Komm mit.“

Ken packte Thorald beim Unterarm und zog den protestierenden König aus seinen Gemächern heraus.

„Bei Mutter Natur, habe ich nicht... nicht soeben gesagt, dass du mich nicht kontrollieren sollst“, murrte Thorald.

„Ja, das hast du“, räumte Ken ein. Mit einem schiefen Grinsen fügte er dann hinzu: „Aber ganz ehrlich, willst du lieber hier drinnen reihern?“\bigskip







Ken führte den König aus Brandurs Turm hinaus, mitten auf die Brücke zwischen dem Turm und dem Thronsaal.

Thorald kniete nieder und blickte über die Brüstung, sein Magen rumorend. Dann spuckte er. Ken stand neben ihm, die Hand auf die Schulter, aber abgewandt. Da kauerte er, der König neben seinem Ratgeber, sich über die Brüstung gebeugt von einem der höchsten Punkte des Reichs in die Tiefe übergebend. Der Mond stand immer noch hoch am Himmel. Tief unter den beiden lag das Dorf der Rietburg still da. Auch das rhythmische Hämmern in Warguths Schmiede hatte aufgehört.

„Besser raus als rein“, warf Ken hilfreich ein, „Bist du endlich fertig?“

Thorald stöhnte schwach. Ken klopfte ihm ungeduldig auf die Schulter und zog ihn dann in die Höhe.

Seine Worte waren hart, doch seine Stimme war urplötzlich wieder weich geworden: „Ich kann dieses jämmerliche Getue nicht länger mit ansehen. Ich weiß, was du brauchst. Komm mit. Na, komm schon, hat dich mein Rat je in Schwierigkeiten gebracht?“

Ken schleifte den Thronerben mehr von der Brücke, als dass er ihn führte. Mit einem eleganten Schlüssel öffnete Ken die Tür zum Thronsaal und führte Thorald hinein. Es war stockfinster, aber Ken schien sich gut auszukennen, denn er verschwand schnurstracks im Dunkeln. Thorald ließ erneut die Schwerkraft auf sich wirken und sank an einer Säule neben dem Eingang zu Boden. Es spuckte aus, aber der säuerliche Geschmack in seinem Mund wurde nicht weniger. Im Schatten fiel es wenigstens nicht auf, wie sehr die Welt wackelte.

Ein Knirschen und Knarren weiter hinten im Thronsaal ließ Thorald aufhorchen.

Bläuliches Licht flammte auf und warf den Schatten einer großen schwarzen Silhouette bis vor seine Füße. Thorald brauchte einige Momente, um zu verstehen, was er sah: Ken hatte die versteckte Klappe hinter dem Thron geöffnet und wuchtete nun den strahlend hell scheinenden Sternenschild dahinter hervor.

„Jetzt komm schon her“, ächzte der Kahlköpfige, „Der Schild scheint sogar zu spüren, dass du ihn jetzt brauchst.“

Thorald schüttelte bloß den Kopf: „N... Nein, Ken, das mache ich nicht noch einmal durch. Mehrere Monate habe ich jetzt ohne den Sternenschild durchgehalten... er kann regelrecht süchtig machen.“

„Hör für einmal auf, dich zu sorgen, und greife nach der Hoffnung, du tumber Troll! Du willst sie doch!“

Thorald schüttelte weiterhin den Kopf, doch das Licht des Sternenschilds zog ihn schon beinahe magisch an. Ehe er sich versah, stand er neben der verborgenen Klappe und blickte auf den Bruderschild hinunter, in dessen verdreckter Verzierung sich der mächtige Schild der Hoffnung widerspiegelte. Und was für ein Licht dieser aussandte! Thorald musste seine Augen zukneifen, um nicht zu erblinden. Fast wie im Traum streckte er seine Hände aus und berührte das sternförmige Emblem des Schilds.

Thorald hatte einen schwachen Anstieg von Stolz in seiner Brust erwartet, wie es die letzten Male gewesen war, bei denen er sich auf den Schild gestürzt hatte. Doch dem war nicht so. Kaum hatten Thoralds Fingerspitzen die Schildoberfläche berührt, glühte der Schild noch heller auf und Thorald fühlte, wie ein riesiges Gewicht von seinen Schultern genommen wurde.

Es war wieder wie beim ersten Mal, als er vorsichtig den Sternenschild berührt hatte. Seine Kopfschmerzen und seine Übelkeit waren immer noch da, aber völlig in den Hintergrund getreten. Sie waren nicht wichtig. Nichts war wichtig außer diesem einen Moment, in dem Thorald den Sternenschild hielt und mit zusammengekniffenen Augen in dessen grell leuchtende Oberfläche starrte. Ein leises Summen war zu hören, mit einer wunderschönen Melodie, welche selbst Grenolins klare, reine Stimme übertraf. Thorald glaubte, im Schild Schemen zu erkennen. Oder wurden diese Bilder vor seinem inneren Auge abgespielt?

Er sah eine Masse von Menschen, welche sich wie ein riesiger brauner Wurm über das Rietland auf die Rietburg zubewegten. Freudenschreie und Marschlieder drangen schwach an sein Ohr. Und angeführt wurde die Prozession von zwei in grüne und blaue Gewänder gekleideten Gestalten, die gemeinsam auf einem edlen weißen Pferd ritten.

Je mehr sich Thorald auf diese Personen konzentrierte, desto klarer wurde das Bild. Schwarzes Haar, ein Bogen an ihrer Seite, und war das etwa ein schwarzer Hund, welcher neben dem Pferd her trottete? Ganz eindeutig, das war Chada, die Heldin von Andor. Zeigte der Sternenschild ihm die Zukunft? Würden die Helden von Andor nach Andor zurückkehren? Aber warum aus dem Süden? Und war das etwa... ja, das war unverwechselbar die goldene Rietgraskrone, welche auf dem Kopf Chadas thronte!

Natürlich! Thorald hatte so lange versucht, aus dem Schatten der Helden von Andor zu treten und Brandurs Fußstapfen auszufüllen, dass er ihr Talent, mit Schwierigkeiten umzugehen, ganz vergessen hatte. An sie könnte er das Königsamt Andors guten Gewissens abtreten. Dann wäre dieser Albtraum endlich vorüber.

Jawohl, dass Chada in dieser Vision die Rietgraskrone trug, war ein eindeutiges Zeichen. Thorald würde sein Amt abtreten und die Krone Andors an die Fürsten Andors abgeben. Chada würde eine würdige Nachfolgerin für den Thron sein. Trotzdem war es bemerkenswert, dass Chada und nicht etwa der gute Thorn, den Thorald schon viel länger kannte und mochte, die Krone trug. Thorald dachte zurück an die Zeit, als er als junger Erwachsener Chada zum ersten Mal getroffen hatte. Damals hatte er schnurstracks davon geträumt, sie eines Tages zu heiraten. Unrealistische Prinzenträume eines selbstverliebten Angebers und Schürzenjägers, natürlich, aber konnte es sein... dass vielleicht eine klitzekleine Chance bestand?

Das Bild verblasste und das leise Summen des Schilds ließ gemeinsam mit dem Glühen nach. Der Sternenschild hatte fürs erste seine Kraft verbraucht. Aber er hatte seine Wirkung getan. Die Hoffnung, die Thorald verspürt hatte, ließ nicht nach.

Die Worte Rekas hallten in seinem Kopf nach: „Du musst vorsichtig sein, Thorald. Ich habe es im Traum erblickt. Eine dunkle Gestalt dürstet nach deiner Thron, und sie wird dir eher Schaden zufügen, als dass sie die Rietgraskrone auf einem anderen Kopf sieht.“

Damals hatte Thorald Rekas Worte als abergläubiges Geschwätz abgetan, aber heute, hier und jetzt, in einer für ihn so unüblichen Klarheit, die der Sternenschild ihm geschenkt hatte, war er sich nicht mehr so sicher. Jetzt, wo er so darüber nachdachte, entstammten nicht nur viele seiner guten Entscheidungen Kens Rat, sondern auch seine schlechteren. Die Helden von Andor hatte er auf Kens Rat in den Norden geschickt. Hatte Ken Thoralds Ansehen zu mehren versucht oder seine stärksten Beschützer außer Reichweite bringen wollen? Und wie oft hatte Ken ihn bereits zu einem Umtrunk angeregt, nur um ihm danach eine Standpauke über dieses Laster zu halten? Wie oft hatte er ihm bereits den Sternenschild überreicht und in den siebten Himmel gelobt? Thorald erschauerte. Sein Selbstbild war wie Wachs in Kens Händen. Seine Entscheidungen auch. Das konnte nicht gut sein. Warum erkannte er das erst jetzt?!

Als das Leuchten des Sternenschilds vollständig erlosch, blieb der Thronsaal kurzzeitig in beinahe vollständige Dunkelheit getaucht und nur das Nachbild des leuchtenden Schildes war erkennbar. Dann flammte wie aus dem Nichts eine Flamme auf und beleuchtete die tiefen Furchen im Gesicht Kens und die Narbe an seiner rechten Wange. Ken hängte die soeben angezündete Fackel in eine Halterung an einer Säule, schnappte Thorald den inzwischen erkalteten Sternenschild aus den Händen und bugsierte ihn vorsichtig in die versteckte Klappe hinter dem Thron.

„Flausen ausgetrieben?“, sprach er Thorald an, den Rücken immer noch zugewendet, „Siehst du nun wieder, dass der Thron Andors dir zusteht und dass noch Hoffnung für das Reich besteht? Dass eine Abdankung nicht nötig ist?“

Thorald nickte schicksalsergeben und Ken klopfte ihm aufmunternd auf die Schulter. Thorald zuckte vor seiner Berührung zusammen.

„Lass...“, krächzte Thorald, „kannst du... ich glaube, ich brauche etwas Zeit für mich selbst.“

Kens Augenbrauen zogen sich zusammen, aber er nickte: „Natürlich. Es ist auch wahrlich an der Zeit, dass ich mir endlich eine Mütze Schlaf gönnen kann.“

Er verbeugte sich und huschte geschwind aus dem Raum. Thorald blieb alleine zurück. Alleine mit seinen Gedanken.

Die Vision des Sternenschilds zeigte ihm deutlich, dass Andor eine gute Zukunft bevorstehen konnte, auch oder gerade ohne Thoralds Führung. Die Krone Andors musste nicht für immer auf seinem Haupt bleiben. Er würde sie an einen würdigeren Nachfolger abtreten, sobald die Helden von Andor aus dem Norden zurückkommen würden. Allzu lange konnte es ja nicht mehr dauern. Er war erst der zweite König dieses jungen Landes, niemand konnte ihm vorschreiben, dass er das Königsamt nicht übergeben könnte. Und dann würde er keinen Schaden mehr anrichten können. Müsste keine langen Sitzungen mit den Anführern der Armee und den Vertretern des Volkes haben. Hätte einfach nur Zeit für sich selbst. Würde er arbeiten müssen? Vielleicht könnte er bei der Ausbildung der Krieger helfen. Es würde bestimmt ein Plätzchen für ihn geben.

Nur sollte er vielleicht lieber vorsichtig sein und Ken nicht darauf aufmerksam machen, was er vorhatte. Ein Trunkener Tor war er, warum nur hatte er Ken von Rekas Warnung erzählt?

Thorald wankte die Säulen des Thronsaals entlang. Da, an der Wand, hing ein Gemälde seiner Eltern. König und Königin, Hand in Hand. Er war aus Krahd geflohen und sie stammte aus den Flusslanden. Sie war friedlich die Narne hinuntergegangen, als Thorald noch nicht einmal ein Schwert hatte führen können. Brandurs Tod war ihm hingegen traumatisch tief in seine Erinnerungen gebrannt. Thorald spürte erneut die Panik, die er ihn damals ergriffen hatte. Späher des König hatten berichtet, wie sich in den Ausläufern des Grauen Gebirges ein Drache aus der Knochengrube erhob. Es war erst Stunden her gewesen, dass sein Vater gefallen war. Der Kampf zur Befreiung der Rietburg war noch vollends im Gange gewesen. Damals hatte er gedacht, dass Tarok sein Ende sein würde. Nur der tapfere Einsatz der todesmutigen Helden hatte ihn gerettet. Wäre es besser gewesen, wenn er damals umgekommen wäre? Wie hatten seine Eltern es bloß geschafft, ein Königreich zu führen, ohne dabei den Verstand zu verlieren? So viele unterschiedliche Wünsche anzuhören und gerechte Entscheidungen zu treffen?

Als Thoralds Blick auf den jungen schwarzhaarigen Burschen auf dem Gemälde fiel, welcher zwischen seinen beiden Eltern stand und stolz in die Welt hinausschaute, da kamen sie endlich, die Tränen. Und wie sie kamen. Als wäre ein riesiger Staudamm gebrochen, so wie Thoralds naiver Stolz schon lange gebrochen war. Der rote Mantel des Königs, den er mit Wein, Met und Gorblut besudelt hatte, lag schwer auf ihm. Brandurs Fußstapfen waren zu groß für ihn gewesen, nicht als ein blasser Schatten seines Vaters war er. Und die Stärke seiner Regierung war nun abhängig von einem Berater, welcher ihm vielleicht lieber die Krone mit Gewalt abnehmen würde, als seinen allumfassenden Einfluss über ihn schwinden zu sehen. Es war zum Verzweifeln.

Thorald stolperte zur Schatulle, in der die Rietgraskrone für zeremonielle Anlässe verstaut war. Er klappte den Deckel hoch und betrachtete die goldene Krone, deren dünne Zacken im Fackellicht umherzuwiegen schienen wie das Rietgras im Wind.

Ein Bild erschien vor seinem inneren Auge: Sein Vater, wie er mit der Krone schief auf dem Kopf und einem Wanderstock in der Hand durch das Rietgras strich. Im Jahre seines Todes hatte Brandur noch versucht, die Wunden der Vergangenheit zu heilen und Frieden mit Fürst Hallgard von den Schildzwergen zu schließen. Ein Frieden, der jahrzehntlang unmöglich erschienen war, weil Brandur den Schildzwergen den Sternenschild einst vorenthalten und dann diesen an den Verräter Rudnar verloren hatte. Kram hatte sich nach Brandurs Eskorte zur Mine dafür eingesetzt, dass der Bruderschild bei den Helden von Andor bleiben konnte, statt in einer Schatzkammer der Schildzwerge zu verstauben. Und nun verstaubten diese beiden mächtigen Schilde in einer Klappe hinter dem hölzernen Thron Andors. Weil Thorald ihren Einsatz verwehrt hatte, als Kram beim Kampf um Cavern danach verlangt hatte. Aus Furcht davor, betrogen zu werden, und weil er den Sternenschild in seinen Zeiten der Hoffnungslosigkeit für sich selbst gebraucht hatte.

Aufgrund dieser Entscheidung waren vermutlich viele Zwerge gefallen. Mit dem Bruderschild hätte man Kram aus weiter Ferne Stärke schenken können, bis die restlichen Helden in Cavern eingetroffen wären. Vielleicht hätte man mit Sternenschild gar Hallgards Tod abwenden können. Nicht, dass Thorald per se viel an den goldgierigen Zwergen läge. Aber sie waren praktische Handelspartner und Thorald hatte es wieder einmal geschafft, mit seinen Entscheidungen die Beziehungen und somit das Schicksal Andors zum Schlechteren zu neigen. Er fluchte und wischte sich das tränennasse Gesicht ab.

Kram war inzwischen nicht nur ein Fürst von Andor, sondern auch der Fürst von Cavern geworden, das Oberhaupt aller Schildzwerge. Aber wenn Thorald seinen Informanten vertrauen konnte, so hatte Krams Ernennung nicht die Anerkennung aller Schildzwerge genossen. Eine Splittergruppe, die sich selbst die „Wahren Schildzwerge“ nannten, hatte schon einmal versucht, ihn zu stürzen, und wer vermochte zu sagen, ob sie es nicht ein zweites Mal versuchen würden? Ihr Anführer war Gerüchten zufolge der stärkste Zwerg, den die Minen je gesehen hatten, und er hasste Menschen aus ganzer Seele. Zumindest hatten Kens Kontakte das behauptet.

Thorald fühlte erneut den Stich des Schwarzen Ritters unter seinem Herzen. Kram war es gewesen, welche die lebensrettenden Zauberhutpilze für Thorald beschafft hatte. Und Kram war es gewesen, der sich mit einem Schild todesmutig zwischen einen Feuerstrahl Taroks und den am Boden liegenden Prinzen geworfen hatte. Kram hatte ihm gleich mehrere Male das Leben gerettet. Thorald schuldete ihm so vieles. Hatte er da nicht eine gewisse Verpflichtung, Kram zu unterstützen? Auch wenn viele Zwerge sich nicht um das Schicksal der Menschen zu kümmern schienen, so war Kram derjenige, der es doch tat. Das Reich konnte nur davon profitieren, wenn Thorald Kram als Fürsten Caverns unterstützte.

Thorald richtete sich abrupt auf und setzte die Rietgraskrone auf seinen Kopf. Ihm war soeben eine geniale Idee gekommen. Brandur hatte es schon einmal geschafft, mit einem Marsch zur Mine einen ihm abgeneigten Fürsten der Schildzwerge friedlich zu stimmen. Thorald konnte es ihm nachtun. Er hatte bereits mehr als bewiesen, sich selbst verteidigen zu können. Und im Gegensatz zu Brandur damals besaß Thorald jetzt sogar etwas, das die Schildzwerge von ganzem Herzen wollten.

Sein Blick schweifte zurück zur geheimen Klappe hinter dem Thorn. War es das wirklich wert, die beiden mächtigen Schilde aufzugeben?

Ja, beschloss Thorald, das war es. Die Beziehung Andors zu Cavern brachte für beide Reiche nur Vorteile, und Kram brauchte Symbole wie die mächtigen Schilde, um die Anerkennung im gesamten Zwergenvolk zu erhalten. Ganz nebenbei war es ein Akt der Rebellion gegen Ken. Ken hatte immer darauf beharrt, dass die mächtigen Schilde in der Rietburg bleiben sollten. Thorald würde ihm hierdurch zeigen, dass immer noch er der König dieses Königreichs war.

So viele Entscheidungen Thoralds hatten schlechte Nachwirkungen nach sich gezogen. Dies war seine Chance, dem Reich, seinem Vater und sich selbst zu beweisen, dass er nicht bloß ein Taugenichts war. Er war der König von Andor, und er würde noch in dieser Nacht nach Cavern aufbrechen. Mit den beiden mächtigen Schilden im Gepäck und der Rietgraskrone auf dem Kopf. Alleine, wie sein Vater vor ihm. Ein Friedensangebot und ein Aufstand gegen Kens Kontrolle.

Die mächtigen Schilde würden ein Geschenk als Wiedergutmachung der Beziehung ihrer beiden Völker sein. Kram würde sie mit Freuden annehmen. Die restlichen Helden würden bald darauf aus dem Norden zurückkehren, auf dass man ihnen die Rietgraskrone überlassen könnte.

Alles würde gut kommen.








\newpage
\section{Thorald kommt vom Weg ab}


Thorald schreckte aus unruhigen Träumen hoch. Das war eigenartig. Eigentlich hatte er schon seit Jahren keine Albträume mehr gehabt. Nicht, seitdem diese gruselige, schattenhafte Gestalt mit den weißen Augen, dem langen weißen Haar und den zwei langen Schwertern ihn endlich nicht mehr Nacht um Nacht im Traume verfolgt und gefragt hatte, wer der Erbe seines Vaters Willens sei. Und diese Gestalt war Thorald inzwischen schon seit über einem Jahr nicht mehr erschienen.

Thoralds unruhige Träume waren nicht das einzige Eigenartige an seinem Zustand. So bemerkte Thorald rasch, dass er sich nicht in seinem gemütlichen Bett in der Rietburg befand, sondern in einer von Morgentau feuchten Böschung am Ufer der Narne lag. Seine Füße waren so kalt, weil sie bereits zum Teil vom Narnenwasser überflutet wurden.

Rasch richtete sich Thorald auf. Die Erinnerungen von letzter Nacht brachen über ihn herein ein und dröhnten seinen Kopf mit Schmerzen zu. Ein Blick auf seine Umgebung bestätigte es: Er war tatsächlich mitten in der Nacht nach Cavern aufgebrochen. Der Bruderschild und der Sternenschild lagen in einem durchnässten Sack an der Böschung, einige Schritte von ihm entfernt. Er hatte es offenbar irgendwie bis über die Markbrücke geschafft, ehe er hier zusammengebrochen war.

Der König räkelte sich und versuchte anhand des Sonnenstands zu erkennen, wie spät es war. Keine Chance. Zu seiner Linken sah er die mächtige Zwergeneiche hoch in den Himmel ragen, zu seiner Rechten erkannte er eine blau gewandete Frau, welche am anderen Flussufer saß und ihn interessiert zu mustern schien. Thorald langte sich instinktiv an den Kopf, und zog scharf Luft ein. Die Rietgraskrone saß nicht mehr da. Er musste sie gestern verloren haben. Ach, warum war er nur zu dieser wahnwitzigen Tat aufgebrochen? Ein einfacher Bote oder sogar ein Falke mit einer Einladung zu Schildübergabverhandlungen hätte doch gereicht.

Aber gut, jetzt, wo er schon einmal hier war, konnte er die Mission auch rietgraskronenfrei zu Ende führen. Thorald kämpfte sich die matschige Böschung hoch und spürte die warme Sonne auf seinem Gesicht. Von hier aus waren es nur noch wenige Stunden bis zum südlichen Eingang Caverns und der Weg war gut ausgebaut.

Es wäre so einfach, die beiden Schilde jetzt Kram abzuliefern.

So einfach.

Aber nun, ein wenig nüchterner, war sich Thorald plötzlich nicht mehr sicher, ob er sie wirklich abgeben wollte.\bigskip







Thorald wandte sich vom Weg nach Cavern ab und folgte dem geschlängelten Flusslauf der Narne gen Norden. In einem Büschel Gras am Ufer fand er einen hübschen grünen Runenstein, der in Thoralds Faust leise vor sich hin summte. Thorald steckte ihn als gutes Omen ein und machte sich auf den weiteren Weg.

Ein Trampelpfad führte von hier aus über die Bogenbrücke in den Wachsamen Wald, doch war dieser nicht Thoralds Ziel. Thorald selbst war sich seines Ziels nicht ganz bewusst, ehe er in der Ferne eine ärmliche Bauernkate am Fuße des Grauen Gebirges erkannte.

Elgas Hof.

Thorald schund sich innerlich. Warum hatte er sich nicht bei ihr gemeldet? Wie viel Zeit war seit ihrem letzten Treffen vergangen? Thorald hatte Elga seit ihrer Heirat immer seltener besucht, und doch erinnerte er sich nun, als wäre es erst gestern gewesen, an eine Zeit zurück, in der er so viel Zeit wie nur möglich bei ihr verbracht hatte. Elga hatte ein hartes Leben gehabt, alleine ihren Bauernhof bestellend, und sie hatte in Thorald nie bloß einen verwöhnten reichen Prinzen gesehen. Ihre gemeinsame Zeit war viel mehr als nur ein Techtelmechtel gewesen, sondern sehr wertvoll für Thorald als Person. Als er nach Brandurs Tod mit der Krone gehadert hatte, war Elga es gewesen, die ihn davon überzeugt hatte, das Amt zu übernehmen. Als die Frau des alten Erwan von gegenüber verstorben war, hatte Thorald eine prunkvolle Beerdigung veranstalten lassen, wie es das östliche Rietland noch nie gesehen hatte. Dann aber hatten die Pflichten eines Königs Thorald mehr und mehr in die Rietburg beordert. Ken hatte mehr und mehr seiner Zeit in Anspruch genommen. Elga hatte ihre eigene Familie gegründet. Und so waren sie auseinandergetrieben worden. Wie es ihr wohl ergangen war in all dieser Zeit?

Kurz kicherte Thorald darüber, wie absurd die Situation war. Der König von Andor klopfte mit Kopfschmerzen bei einer einfachen Bäuerin an der Tür, mit den zwei mächtigsten Artefakten des Landes bei sich und höchstwahrscheinlich bald einem verärgerten Suchtrupp der Rietgarde auf den Fersen. So nahm er seinen Mut zusammen und pochte an die verfallene Holztür.

Es dauerte nicht lange, bis die Tür aufflog und Thorald in zwei blitzende braune Augen starrte, die definitiv nicht zu Elga gehörten. Die Frau blickte kurz überrascht drein, als sie Thorald erkannte, fasste sich allerdings ebenso rasch und deutete dann einen Knicks an: „Euer Hoheit!“

Thorald fragte sich, wie viel Ironie hinter der Höflichkeitsgeste steckte.

„Juna! Es ist schön, dich... ich... ist Elga hier?“, fragte Thorald angespannt.

Juna legte ihren Kopf schief und meinte fröhlich: „Du wirst lachen, aber ich bin mir nicht einmal sicher.“

Sie drehte sich um und rief laut ins Innere der Kate: „Elgalein! Bist du drinnen? Dein König ist hier!“

Es blieb einen Moment lang still. Juna wandte sich bereits wieder Thorald zu und setzte zu einer Absage an, als doch noch Schritte ertönten und die Tür zur Gänze aufgerissen wurde. Da stand Elga und lächelte Thorald breit an. Thorald versuchte, ein Lächeln aufzusetzen, aber es wollte ihm nicht vollends gelingen.

Elga winkte ihn dennoch hinein.\bigskip







So saßen sie bald beide an Elgas altem Holztisch, an dem sie schon so oft gesessen hatten. Thorald hatte die mächtigen Schilde und sein Schwert zur Seite gelegt, konnte seine Hände nun nicht still halten und kratzte an der rauen Tischplatte herum. Elga stand still da und musterte ihn aufmerksam, aber ruhig. Juna tigerte hinter Elga umher und blickte immer wieder verstohlen zu Thorald hinüber.

Es war Elga, die das Wort ergriff.

„Ach, Thorald, du hast schon besser ausgesehen. Was ist mit dir geschehen?“

Thorald fragte sich, ob diese Worte sich auf seinen jetzigen verdreckten Zustand bezogen und die zwei Schilde im großen Sack neben seinem Holzstuhl, oder darauf, wie er Elga mit der Zeit weniger und weniger besucht hatte. Er beschloss, sich erst einmal um ersteres zu kümmern, und antwortete knapp.

„Ich war auf dem Weg nach Cavern. Ich hatte gestern betrunken die stolze Idee, die mächtigen Schilde ihren... ‚rechtmäßigen‘ Besitzern zurückzubringen.“

Elga nickte bloß. Juna hatte aufgehört, hin- und herzuwandern, und starrte Thorald nun direkt an. Thorald warf ihr einen knappen Blick zu und fuhr dann fort, versuchte, die passenden Worte zu finden.

„Die Wahrheit ist... ich... ich weiß nicht, was ich tun soll. Ich habe nicht viele gute Entscheidungen getroffen, und die meisten guten verdanke ich nur meinem Berater. Es steht meinetwegen nicht gut um die Zukunft des Landes und die Helden von Andor sind wegen mir weit weg. Ich dachte, dass das Volk mich dadurch mehr ehren würde, aber dieses wendet sich immer mehr von mir ab. Ich halte das nicht mehr aus.“

Er beugte sich vor und atmete tief durch. Elga flüsterte etwas zu Juna, welche sich umdrehte und langsam aus dem Raum lief, aber dabei immer wieder mürrisch nach hinten sah. Offenbar hätte es sie brennend interessiert, wie Thoralds Geschichte weiterging.

„Ich glaube... ich glaube, dass die Krone Andors auf einem anderen Haupt besser sitzen könnte als auf meinem Brummschädel“, brummte Thorald dann, „und das vergesse ich nicht einmal im größten Trinkgelage. Ich... ich will sie abtreten. Aber das kann ich erst, wenn die Helden zurückkehren. Und... ich beginne, mich vor Ken zu fürchten. Ich wüsste nicht, an wen ich mich wenden sollte, wenn er nicht an meiner Seite wäre und mich beraten würde, aber ich werde das Gefühl nicht los, dass er etwas... er hat eine finstere Seite, eine manipulierende, und diese kommt immer öfter zu Vorschein.“

Juna betrat den Raum erneut. Sie trug eine Platte mit Milchkrügen und Käsescheiben bei sich. Thorald spürte erst jetzt, wie trocken seine Kehle war und wie sehr sein Magen knurrte. Juna vollführte erneut einen kleinen Knicks und bot ihm die Fressplatte an. Thorald langte zu und bedankte sich überschwänglich.

Elga sprach, und ihre klangvolle Stimme beruhigte Thorald etwas: „Das Los der Krone hast du noch nie leicht getragen, auch wenn du das früher gerne geglaubt hast. Bei all den Trinkereien, in die du stürzt, bei dem makellosen Ruf deines Vorgängers, bei deinen Ambitionen und bei deinem listigen Berater, da ist es wahrlich kein Wunder, dass es dir schwer fällt, dich so um dein Land zu kümmern, wie du es dir gerne wünschtest... oder deine Beziehungen.“

Thorald blickte beschämt auf, doch konnte er in Elgas blitzenden Augen keine Wut erkennen, bloß einen Hauch von Traurigkeit. Noch etwas, was sie so stark von Ken unterschied.

„Es tut mir leid, dass ich mich nicht mehr bei euch gemeldet... Wie... wie geht es euch? Wie ist es euch ergangen?“

„Wir kommen klar“, meinte Juna mit einem aufrichtigen Lächeln, griff sich einen Batzen Käse und zog sich dann langsam wieder aus dem Raum zurück.

Elga bestätigte: „Wir kommen klar.“

Dann sanken ihre Mundwinkel: „Gerüchte von verschwundenen Bauern im östlichen Rietland und von wilden Wölfen an den Ufern der Narne gehen um. Der alte Erwan von nebenan ist gerade verstorben. Man munkelt, die Wölfe hätten ihn gerissen. Aber seine Kinder sind stark. Wir gucken immer wieder bei ihnen vorbei. Sie scheinen ebenfalls klarzukommen.“

„Das... das tut mir leid“, brachte Thorald hervor, „Ich hätte hier sein sollen. Kann ich... lässt sich eine Beerdigung...“

„So sehr uns das Angebot doch ehrt, so glaube ich nicht, dass eine weitere prunkvolle Beerdigung das ist, was seine Kinder sich wünschen. Ganz abgesehen davon hat euer Wolfskrieger schon eine schöne Zeremonie abgehalten.“

Thorald nickte überrascht. Orfen konnte er sich schlecht bei einem Beerdigungsritual an der Narne vorstellen – aber hinter dessen rauer Fassade versteckten sich offenbar so einige verborgenen Talente.

„Dann... braucht ihr... kann ich euch sonst irgendwie helfen? Ich könnte eine Patrouille der Rietgarde hier vorbeischicken, um sich diese Wölfe einmal anzuschauen.“

„Danke, aber ich glaube nicht, dass das nötig ist. Juna meistert den Umgang mit ihrem Stab immer besser. Bald wird sie soweit sein, die Wölfe friedlich zurückzutreiben – sollten diese sich überhaupt je einmal mit Hintergedanken auf unseren Hof wagen.“

Hatte Thorald wieder einmal versucht, mit Geschenken und Gaben sein schlechtes Gewissen zu besänftigen? Elga hatte ihn schon einmal auf diese Tendenz angesprochen, lang, lang war’s her.

Eine kurze Zeit lang blieb es still, während Thorald sinnierte. Dann fuhr Elga fort: „Was du tun könntest, ist, wieder öfters hier aufzukreuzen. Das Leben eines Königs besteht natürlich nicht aus sehr viel Freizeit, aber es würde uns freuen, dich wieder einmal hier begrüßen zu dürfen, statt dass du Gildas Met oder gar erneut dem Starkbier der Zwerge verfällst. Dann hättest du auch etwas Zeit weg von diesem Ken. Nach den wenigen Gerüchten, die uns von der Rietburg zufliegen, hat er dich ziemlich im Griff. Wie man sich doch in einem Menschen täusche kann. Ich mag mich noch gut an die Zeit erinnern, als du ihn kennenlerntest und stundenlang von ihm schwärmen konntest. Damals erschien er noch wie der freundliche junge Bursche, der gerade zum richtigen Zeitpunkt zu dir gefunden hatte. Er beriet dich gut und verschaffte dir kleine Siege, die dein Selbstvertrauen stärkten. Er schaffte es, dass dir die Fußstapfen deines Vaters nicht so riesig vorkamen, als dass du darin versinken müsstest.“

Thorald schluckte. An die schöne Zeit, als er Ken kennengelernt hatte, wollte er nicht zurückdenken. Zu kompliziert war sein Verhältnis mit Ken inzwischen geworden. So wandte er sich von diesem Thema und sprach stattdessen Elgas Vorschlag an, sich wieder häufiger bei ihnen blicken zu lassen: „Das... nein, ihr könnt mich nicht einfach wieder in eure Leben lassen... du sollst mir nicht einfach verzeihen. Ich habe schlimme Fehler begangen und dich und Juna im Stich gelassen. Ihr solltet wütend sein auf mich... alle anderen sind’s doch auch...“

Elga lächelte traurig: „Ach, Thorald... keiner leugnet, was für Schäden du in deinem Leben angerichtet haben magst. Aber du bist immer noch eine wertvolle Person mit einem guten Kern. Dich in ein Loch der Schuldgefühle reinzuwerfen, könnte vielleicht einigen Genugtuung verschaffen, wäre aber ziemlich sicher nicht hilfreich für dich oder das Königreich. Vielleicht bist du ja hier, um ein klein wenig Hilfe dafür zu erbitten, aus diesem Loch herauszuklettern? Wir helfen gerne, wo wir können.“

„Ich brauche mehr als nur ein klein wenig Hilfe“, murmelte Thorald, „Und ich verdiene sie ohnehin nicht. Ich will... ich will einfach nur die Krone loswerden. Nun, das hat mein betrunkenes Ich gestern Nacht tatsächlich geschafft. Die Rietgraskrone habe ich auf dem Weg hierher verloren. Aber das Amt abgeben steht noch an... ach Elga, ich weiß einfach nicht weiter.“

Elga stützte sich auf ihr Kinn und murmelte: „Hui, du steckst wirklich in einem unschönen Mentalität fest.“

Thorald klaubte weiter an der hölzernen Tischplatte herum und wich Elgas Blick aus, als diese fortfuhr: „Aber das ist in Ordnung. Jeder braucht Hilfe von Zeit zu Zeit, und dass du hierhergekommen ist, ist schon ein guter Anfang. Kleine Schritte führen zum Ziel, und du bist es wert, diesen Weg zu gehen. Zuerst einmal wollen wir doch herausfinden, was du wirklich mit den mächtigen Schilden anstellen willst, und ob dir Ken Dorrs Nähe mehr bringt oder schadet. Magst du hier bleiben für den Moment? Ich muss diesen Nachmittag noch etwas Saatgut streuen und natürlich die Kühe melken – da könntest du gerne aushelfen. Am Abend könnten wir uns dann die Zeit nehmen, deine Optionen zu besprechen. Schilde abgeben oder behalten. Krone abtreten oder akzeptieren. Ken Dorr bessere Umgangsformen beibringen oder von dir fernhalten. Wir hätten auch ein Gästebett für die Nacht.“

Thorald lehnte sich zurück. Solche Worte hatte er vermisst, ohne sich dessen überhaupt bewusst zu sein. Ein respektvolles Gespräch über seine Entscheidungsmöglichkeiten, was mehr konnte er sich wünschen. Er nickte langsam.

Da stürzte Juna durch die Tür hinein, ein blankes Schwert in der einen und ein knorriger Holzstab in der anderen Hand. Sie sprach aufgebracht: „Ich unterbreche euch zwei Turteltakuri ja nur äußerst ungern, aber da draußen marschiert eine ganze Schar von Pferden an. Ein Glatzkopf führt sie an. Den mögen wir nicht.“

Thorald sprang auf, während Elga ihren Stuhl zur Seite schob und sich sorgsam erhob. Schon war von draußen ein Wiehern zu hören. Schwere Stiefel klatschten auf den Dreck vor der Hütte und Schritte kamen immer näher.

Juna packte ihren Stab fester und flüsterte zu Elga: „Ich könnte es mit einem magischen Blitz versuchen. Gestern konnte ich ihn halbwegs kontrolliert feuern.“

Elga grinste schief und meinte: „Letzten Endes ist das doch ein lebendiger Mensch mit Wünschen und Träumen, mit dem wir es hier zu tun haben. Ein hochrangiger noch dazu. Ein magischer Blitz scheint mir als Begrüßung etwas... übertrieben? Was meinst du, Schatz?“

Juna nickte mit gesenktem Blick und legte ihren Stab enttäuscht zur Seite.

Thorald richtete sich seine Frisur.

Es klopfte.

Wie schon bei Thorald zuvor öffnete Juna die Tür. Thorald erkannt sofort Kens schnarrende Stimme: „Ist er hier? Richtet ihm aus, dass er etwas vergessen hat!“

Ein metallisches Klingen ertönte. Thorald vermutete, dass Ken soeben die Rietgraskrone zu Boden geschleudert hatte. Das konnte heiter werden. Ken war mal wieder in einer üblen Laune.

Die Tür zur Bauernkate wurde aufgerissen und Thorald erkannte die mächtige Silhouette seines engsten Ratgebers, der ins Dunkel der Hütte starrte und blinzelte. Seine Augen mussten sich wohl zunächst noch an das wenige Licht anpassen.

„Thorald? Bist du da?“, sprach er unwirsch.

Thorald räusperte sich und versuchte, so ehrwürdig wie möglich „Ja“ zu antworten. Kens Stirn furchte sich und seine Stimme sprang eine ganze Oktave in die Höhe, als er rief: „Wo hast du Trampeltier die mächtigen Schilde gelassen?! Sag bloß, dass du sie auch noch verloren hast!“

„Sie sind hier“, antwortete Thorald, „Und bitte, lass den Ton stecken. Du musst nicht so...“

„Ich spreche, wie mir der Schnabel gewachsen ist! Komm Thorald, auf zurück in die Rietburg! Zuerst schläfst du dir deinen Rausch aus und dann unterhalten wir uns einmal gründlich darüber, wie viel Schaden...“

„Ne, Ken, ich war gerade dabei...“

„Wie bitte?“

„Nein, Ken, ich komme noch nicht mit. Ich... ich bin immer noch dein König, verflucht noch mal! Und wenn ich die mächtigen Schilde lieber ihren Urhebern schenke, als dass du sie noch länger zu meiner Manipulation nutzt, dann hast du mir da nicht zu widersprechen!“

Jetzt war es raus. Thorald atmete schwer und wurde sich bewusst, dass er sich nun wieder sicher darin war, die mächtigen Schilde an Fürst Kram abtreten zu wollen.

Ken hingegen schüttelte bloß seinen Kopf und murmelte etwas unter seinem Atem. Thorald besah angespannt, wie Ken nach seinem Schwertgriff langte. Überlegte er sich gerade, ob es geschickter wäre, den unterwürfigen Ratgeber zu mimen oder zu mehr Beleidigungen und Gewalt zu greifen? Thorald wollte für keine der beiden Seiten Kens hier blieben.

So sprang Thorald auf, griff sich den Sack mit den beiden mächtigen Schilden – sein Schwert klapperte zu Boden – raste an Ken vorbei und stürzte ins Freie. Seine Augen brauchten einen kurzen Augenblick, um sich zu orientieren. Während das grelle Leuchten des Sonnenlichts abschwoll, erkannte Thorald die Schemen zweier weiterer Männer in der Rüstung der Rietgarde, die sich gerade zwei Pferden zuwandten, sowie ein drittes, etwas abseits stehendes Pferd, welches wohl Ken zur Anreise genutzt hatte.

Thorald mochte nicht mehr der beste Reiter des Landes sein, aber sofern Kens Handlanger zwei durchschnittlichen Gardisten waren, konnte er ihnen locker das Wasser reichen. Zudem kannte er Kens Pferd gut. So hopste er in der allgemeinen Verwirrung rasch auf Kens Rappen zu, schwang sich mehr oder minder elegant auf dessen Sattel und gab Fersengold.

Thorald lachte wie verrückt, als ihm einmal mehr die absurde Lage bewusst wurde, in der er sich befand. Der König des Landes, wie er mit zwei mächtigen Schilden vor seinem Berater davonpreschte. Er wollte einfach nur weg von alledem, doch nicht einmal bei Elga hatte er Ken und der Krone entkommen können. Wohin nun?

Thorald erinnerte sich noch verschwommen an die gute alte Zeit, als sein Vater Brandur mit ihm den Baum der Lieder aufgesucht hatte und Thorald mit einem gewissen Merrik davongeschlichen war, um die Quelle des Likko zu untersuchen. Sie hatten dort eine geräumige Höhle gefunden, welche sich hinter der Quelle verbarg. Merrik hatte die Höhle mit Freude zu skizzieren versucht, während Thorald die Höhlenwände mit kruden Zeichnungen bekritzelt und dann Merrik mit kalten Wasser beworfen hatte. Gute alte Zeiten. Thorald verspürte Nostalgie und hielt die Höhle für einen guten Ort, um sich zu verstecken und kurz zu verschnaufen. Und so tat er es.\bigskip







Thorald hatte ein bisschen Mühe damit, sein Pferd durch den Wasserfall vor dem Höhleneingang zu locken, aber mit ein wenig Zucker und ein wenig Zerren konnte er auch dieses Hindernis überwinden. Ganz nebenbei hatte der Wasserfall den Effekt, Thoralds verdreckte und verschwitzte Kleidung in Wasser zu tränken. Zumindest einen kleinen reinigenden Effekt musste das doch haben.

Im Innern der Höhle angekommen, ließ Thorald sich erst einmal fallen und rieb sich den schmerzenden Schädel. Draußen hatte ihn die warme Nachmittagssonne gewärmt, doch hier drinnen war es nun überraschend kühl. So schlang er seinen Mantel enger und schwang seine Arme, um sich etwas Wärme zu verschaffen. Dann wandte er sich Kens Pferd zu und versuchte, dessen Fell trocken zu rubbeln. Ungeschickt, wie er war, löste sich der Sack mit den zwei Schilden, welcher scheppernd auf den Höhlenboden polterte.

„Beim Barte des Urtrolls!“, fluchte Thorald, aber das laute Rauschen des Wasserfalls versicherte ihm, dass er von draußen nicht gehört werden konnte. Bestimmt war Kens Suchtrupp ihm bereits dicht auf den Fersen. Sollte er bei Melkart Exil suchen? Nein, er war der König von Andor, und er sollte erhobenen Hauptes zur Rietburg zurückkehren und seine Angelegenheiten richten. Noch immer plante er, die Krone an die Helden abzugeben, sobald diese aus dem Norden zurückkehrten. Wie lange das noch dauern konnte... schwer zu sagen. Gildas Gerüchten zufolge hatten sie viel Zeit damit verplempert, verschiedenen Nebelinseln auszuhelfen. Zumindest hatten sie als Belohnung den Sturmschild verliehen bekommen, den dritten mächtigen Schild, das war doch etwas. Ach ja, genau, die beiden mächtigen Schilde hatte er an Kram übergeben wollen, ehe ihn seine Beine zu Elgas und Junas Hof gelenkt hatten. Das Treffen mit Elga war sicherlich kein Fehler gewesen und hatte ihm ein wenig stark benötigte Zuversicht verleihen können, jetzt jedoch war es an der Zeit, sich wieder auf seinen Plan zu konzentrieren.

Ja, er würde die beiden mächtigen Schilde nach Cavern bringen. Sobald es draußen wieder dunkler geworden war. Damit Ken und sein Trupp ihn nicht vorher fanden und davon überzeugten, zurück zur Burg zu kommen und das Vorhaben aufzugeben. Das war ein sinnvoller Plan. Nur dem Warten sah Thorald nicht mit Freude entgegen. Warten war noch nie seine Stärke gewesen. Sein Magen knurrte bereits erneut und sein Blick wanderte über das Höhleninnere, in der Hoffnung, etwas Essbares zu finden. Irgendeine nette Pflanze oder einen Höhlenpilz.

Jetzt erst fiel Thorald auf, wie gut er das Höhleninnere erkennen konnte. Als er als kleines Kind hier gewesen war, war es stockfinster gewesen, aber nun konnte Thorald problemlos die hintere Höhlenwand ausmachen... sowie einen Gang, der davon wegzuführen schien. Aus diesem schien schwach grünliches Licht, als würden sich weiter hinten grün schimmernde Lichtquellen befinden. Er konnte doch unmöglich so nahe am Geheimen See sein, oder?

Ein unterirdischer Gang war nicht zwingend eine gute Nachricht. Höhlengänge bedeuteten potentiell die Präsenz von Höhlenwichten, oder schlimmer gar, Arpachen. Im besten Fall führte der Gang aber über eine Zwergentür direkt ins Herzen von Cavern, und Thorald müsste nicht einmal bis in den Abend warten, um die mächtigen Schilde dort abzuliefern. Er wog kurz seine Optionen ab und beschloss, zumindest einige Minuten in den Höhlengang vorzudringen.

Sein Pferd war rasch wieder mit den beiden mächtigen Schilden bestückt. Der Höhlenboden war relativ flach, somit hatte er kein Problem damit, die Stute in einen erleuchteten Höhlengang zu führen. Sie sträubte sich zunächst, war aber gut dressiert und folgte dann seinem Zügel.

Thorald wusste nicht, wie lange er den Gang entlangstrich. Der grünliche Schimmer wurde manchmal schwächer, dann wieder stärker, und von Zeit zu Zeit glaubte Thorald, Runen zu erkennen, die in die Wände geritzt waren und von denen das schwache grüne Leuchten ausging. Das beruhigte ihn. Demnach waren diese Wände eindeutig von Zwergen bearbeitet worden, wenn nicht sogar von ihnen geschaffen. Er war sicher hier. Es war nur seltsam, dass nirgendwo Fackeln hingen. War dieser Gang etwa so alt, dass die Zwerge beim Graben noch nicht das Geheimnis des Feuers besessen hatten? Thorald versuchte zurückzudenken an die Geschichtslektionen, die ihm einst erteilt worden waren. Schwertkunde hatte er deutlich lieber gelernt gehabt und die Kopfschmerzen vom gestrigen Gelage meldeten sich schon wieder deutlicher zu Wort, also gab er bald auf, die Runen zuordnen zu wollen.

\jdh{1193}

Bei seinem Brummschädel brauchte Thorald auch ziemlich lange, um zu erkennen, dass der Gang sich plötzlich weitete und wieder in den Wald öffnete. Thorald spitzte die Ohren, aber andere Pferde oder auch nur Schritte eines Suchtrupps Kens waren so wenig zu hören, wie sie nicht zu sehen waren. So wagte Thorald es, sein Pferd langsam aus dem Höhlenschatten ins Licht zu führen und tief durchzuatmen. Die Luft im Höhlengang war stickig gewesen, hier draußen im Freien war es viel angenehmer. Es lag allerdings ein ungewohnter süßlicher Duft in der Luft, den Thorald nicht näher zuordnen konnte. Kurz überlegte er sich, zurück in den Höhlengang zu kehren, aber die durch die Blätter knapp erkennbare Sonne warf schon lange Schatten. Bald würde es dunkel sein, da konnte sich Thorald doch schon jetzt auf den Weg nach Cavern machen.

Oder sich zumindest daran machen, seine Orientierung zu finden. Er befand sich irgendwo im Wachsamen Wald, also sollte er, wenn er sich einfach nach Süden bewegte, früher oder später wieder am Likko ankommen. Oder?

War das hier überhaupt der Wachsame Wald? Spuren von Bewahrern sah er nämlich nirgends. Und nach Mammutbäumen sah das Gewächs hier auch nicht aus. Wuchsen im Barbarenland andere Baumsorten? Oder konnte er die ganze Strecke bis zum Südlichen Wald unterirdisch überquert haben? Das schien ihm alles unwahrscheinlich.

Kurz überlegte Thorald, einen Baum zu erklimmen und die Gegend zu überblicken, sah dann aber davon ab, als er sich daran erinnerte, wie solche Kletterpartien in seiner Kindheit üblicherweise geendet hatten. So schwang er sich auf sein Pferd und ritt langsam weiter.

Kaum hatte Thorald ein bisschen Distanz zwischen sich und den Höhlenausgang gebracht, ertönte ein lautes Knirschen. Überraschte blickte Thorald nach hinten und erkannte, dass der Höhleneingang, aus dem er soeben getreten war, Teil eines riesigen steinernen Gesichts war, das in einen Felsen gemeißelt worden war. Thorald war offenbar aus dessen offenem Mund getreten, doch nun schloss sich ebendieser Mund knirschend und verschloss den Höhleneingang.

Erstaunt und erschrocken ließ Thorald sein Pferd näher treten. Sobald er einige Schritte vom Gesicht entfernt war, öffnete sich der Mund des riesigen Gesichts wieder und enthüllte den Höhlengang. Eigenartig. Dieses Mal glaubte Thorald zusätzlich, ein lautes Summen zu vernehmen. Er suchte in seinen Taschen nach dem Urheber des Geräuschs und fand den kleinen grünen Runenstein, den er früher am Narnenufer gefunden hatte und nun vor sich hin vibrierte. Vorsichtig ließ Thorald das Pferd einige Schritte zurückweichen, woraufhin das Summen des Runensteins versiegte und das Steingesicht seinen Mund wieder schloss. Ein faszinierender Mechanismus, um den Eingang zu verstecken, kein Zweifel, aber von den Runenmeistern der Zwerge war auch nicht weniger zu erwarten. Thorald zog an den Zügeln und führte das Pferd vom steinernen Gesicht weg.

Da ertönte ein lautes Plätschern und wie aus dem Nichts ergoss sich ein Schwall Wasser darüber. Ein wahrer Wasserfall verdeckte das steinerne Gesicht mit dem Höhleneingang nun vollständig. Thorald beachtete den Vorgang kaum und versuchte stattdessen orientierungslos, sich für einen optimalen Pfad zu entscheiden. Er wollte schließlich nach Cavern, und diese Mine lag südlich des Wachsamen Waldes. Süden war links von dort, wo sich die Sonne hinbewegte, oder?\bigskip







Thorald erreichte bald die nächste Waldlichtung und musterte interessiert einen blau schimmernden Runenkreis am Waldboden. Er ließ sein Pferd kurz anhalten und versuchte erneut erfolglos, sich zu orientieren. Da trat ein wahrer Hüne von einem Mann auf die Lichtung und bewegte sich schnurstracks auf Thorald zu. Der Hüne trug einen blauen Mantel mit einem über seine Schultern drapierten kurzen braunen Umhang, aber was darauf saß, zog Thoralds Aufmerksamkeit viel eher auf sich. Eine krumme Nase stach aus einem kantigen Gesicht hervor, welches von dichtem braunen Haar umrahmt wurde. Der Hüne stammte in jedem Fall nicht aus der Rietgarde oder zu Ken Dorrs Handlangern. Vielleicht ein Holzfäller?

„Zum Gruße, der Herr“, rief Thorald freundlich, „Mir scheint, ich habe mich etwas verlaufen. Mögt Ihr mir mitteilen, in welcher Richtung der Likko fließt?“

„Schweig stille, Adliger, ‘s gibt kein‘ ‚Likko‘ hier“, brummte der Hüne. Er blieb vor Thoralds Pferd stehen und zog in einer fließenden Bewegung zwei mächtige Schwerter von seinem riesigen Rücken. Inzwischen sah er relativ bedrohlich drein. Als Thorald an seinen Gürtel griff, stellte er erschrocken fest, dass da kein Schwert befestigt war. Er musste es bei seinem hektischen Aufbruch von Elgas und Junas Bauernkate vergessen haben. Trolldreck!

Der Hüne griff nun sanft nach den Zügeln von Thoralds Pferd und verkündigte freundlich: „Falls du’s noch nicht kapiert hast: Das hier is’n Überfall. Ich klau‘ jetzt dein‘ Sack.“

Der Hüne wollte doch tatsächlich den Sack mit dem Sternenschild und dem Bruderschild klauen. Thorald hätte beinahe aufgelacht.

„Weißt du denn nicht, wer ich bin?“, fragte er nervös.

„Irgendein stinkreicher Faulpelz, das biste“, meinte der Hüne selbstsicher und griff frech nach dem Sack mit den mächtigen Schilden.

Jetzt hatte Thorald genug. Geschwind löste er den Sack vom Sattel des Pferds, schwang sich aus selbigem und warf sich einige Schritte vom Hünen entfernt in Pose. Kurz überlegte sich Thorald, davonzurennen, aber dann würde er nur noch verirrter sein. Und wenn er sich diesen kräftigen Hünen ansah, würde Thorald wahrscheinlich die Schilde zurücklassen müssen, um ihm entkommen zu können. Die Schilde auch noch zu verlieren war eine Schande, das Thorald einfach nicht zulassen konnte. Es war schade, dass Thorald sein Schwert bei Elga und Juna vergessen hatte, aber das sollte nicht heißen, dass er nicht einem dahergelaufenen Rabauken das Wasser im Faustkampf reichen konnte.

„Ich bin dein König, du tumber Troll!“, protestierte Thorald, „Und wenn du diese Schilde willst, so musst du durch mich hindurch.“

„Wie du willst“, schmunzelte der Hüne. Er legte seine beiden Schwerter sorgfältig zu Boden („D‘mit ich dich nicht aus Versehen ernsthaft verletz‘. Wir sind kein‘ Bösen, wir woll‘n nur dein Hab und Gut“) und bewegte sich dann überraschend wendig auf Thorald zu.

Thorald hob seine geballten Fäuste und wich vorsichtig zurück.

„‘Der König‘, als ob ich nich‘ lach‘“, murmelte der Hüne und griff nach Thorald.

„Es stimmt!“, erwiderte Thorald trotzig, während er dem Hünen aus dem Griff glitt und ihn mit seinen Fäusten bearbeitete, „Und wenn du mich nicht sofort gehen lässt, wirst du vor den Rat der Bewahrer kommen, dass schwöre ich!“

„Der König ist leider weit weg von hier“, murmelte der Hüne unberührt, „Und wenn die Krone mich fängt, krieg‘ ich ohnehin schon den Strick. Der steht allen Geächteten zu.“

Thorald bemühte sich, die Kniekehle des Hünen zu treffen, aber dieser ließ sich davon nicht groß beeindrucken. Dann fand die Faust des Hünen ihr Ziel und Thorald wurde schwarz vor den Augen.













\newpage
\section{Thorald trifft vom Weg Abgekommene}

„Edler Herr, seid Ihr verletzt?“

Thorald rieb sich zum zweiten Mal an diesem Tag den schmerzenden Kopf und erwachte aus der Dunkelheit. Über ihm gebeugt stand ein Soldat in einer ihm unbekannten Uniform und stupste ihn mit einem Speer an. Rot war seine Kleidung, und rot-weiß war das spitze Schild, das er fest in seiner anderen Hand hielt.

„Wo... wo bin ich?“, stammelte Thorald.

„Im Forest seid Ihr, und allem Anschein nach hat Euch einer der Geächteten übel erwischt. Diese Bastarde haben es schon seit Wochen auf die Adligen der Grafschaft von Nottingham abgesehen.“

Die Wache legte ihren Speer ab und kniete neben Thorald hin: „Geht es Euch gut? Wurdet Ihr am Kopf erwischt?“

Thorald nickte bloß und fragte dann: „Konntet Ihr die magischen Schilde sicherstellen? Und wisst wenigstens Ihr, dass ich Euer König bin?“

Die Wache blickte verdutzt drein und murmelte dann: „Auweia. Öhm. Kommt, Ihr müsst Euch erholen. Ich besitze leider nicht die Heilkünste, die benötigt werden, um mit einer solchen Kopfverletzung umzugehen. Stützt Euch auf mich, ich bringe Euch zur Burg. Unser Heiler wird sich Euer annehmen können.“

„Ja, die Rietburg ist ein gutes Ziel, ich weiß, wo die steht“, lallte Thorald. Sein Blickfeld wurde wieder unscharf, und die Stimme der Wache hallte, als befänden sie sich immer noch in der Höhle bei der Quelle des Likko.

„Ja, ja, zur Burg mit uns. Auf jetzt. Nein, lasst Euch nicht fallen. Bei Gott, Ihr seid ja in einem Zustand...“

Danach nahm Thorald nur noch verschwommen war, wie Sprachfetzen an sein Ohr drangen.

„Warum müsst Ihr auch so schwer...

„Kommt, Alrich, helft mir kurz...“

„Dieser Karren ist konfisziert im Namen des...“

Plötzlich musste Thorald nicht mehr laufen, sondern konnte sich auf ein Strohbett sinken lassen. Aber das war gar kein Strohbett, das schaukelte ja viel zu heftig. Wie eigenartig.

So sank Thorald in einen unruhigen Schlaf.\bigskip







Als Thorald aufwachte, brauchte er einen Moment, um seine Erinnerungen zu sortieren und zu verstehen, was geschehen war. Es gelang ihm nicht vollständig. Er bemerkte, dass sein Kopf nicht mehr schmerzte und stattdessen von einem weichen Samtkissen aufrecht gehalten wurde.

Überrascht schlug Thorald seinen Augen auf und blinzelte gegen die grellen Lichtstrahlen. Dann fiel ihm auf, dass er sich in einem ihm vollkommen unbekannten Raum befand. Und dass eine ihm vollkommen unbekannte Person auf einem Stuhl neben seinem Bett saß und ihn aufmerksam musterte. Der Mann trug einen eleganten, königlich roten Mantel und feuerrotes Haar, das Thorald an Fenn erinnerte. Seine Hände hatte er gefaltet und seine Aufmerksamkeit galt ganz dem erwachenden Thorald.

Thorald beschloss, die unangenehme Stille zu unterbrechen. Er räusperte sich, hustete und stammelte dann: „Ich... ich danke Euch ganz herzlich für die Gastfreundschaft. Ich war nur auf der Durchreise durch den Wald, als ich von einem groben Hünen überfallen wurde.“

Die Person auf dem Stuhl nickte: „Das wurde mir von der Wache auch berichtet. Ich bin der Sheriff von Nottingham, und es ist meine Aufgabe, diese Geächteten und ihren Anführer Robin Hood zu fangen und zu richten.“

„Sehr erfreut, Sheriff“, nickte Thorald, „Ich bin der Thorald von Andor“. Ein seltsamer Name war das, Sheriff. Und auch von einem Nottingham hatte er noch nie gehört. Konnte es sein... vor vielen Jahren hatten die Andori erfahren, dass das Fahle Gebirge ihnen ein ihnen bislang unbekanntes Land verborgen hatte, das mythische Tulgor, Land der Architektur und Technologie, der Temm und der Takuri. Konnte es sein, dass die Andori von der Existenz eines weiteren Landes nicht wussten? Auf jeden Fall lag Thorald auf einem Bett in einem ihm fremden Gemäuer. Einer Burg? Thorald wusste von keiner bewohnten Burg außer der Rietburg. Wie fern von seinem Heimatland befand er sich nur?

„Thorald von Andor, soso“, murmelte Sheriff, „Ihr seht eindeutig wie ein Adliger aus, wenn auch ein ziemlich verdreckter. Schon alleine Euer Umhang dürfte mehr wert sein, als eine Wache in einem Monat verdient. Seltsam ist bloß, dass ich noch nie von einem Adelsgeschlecht namens ‚Andor‘ gehört habe. Und ich rühme mich, in den Adelsgeschlechtern ziemlich bewandert zu sein.“

Thorald wusste, dass er vorsichtig vorschreiten musste. Sheriff schien eine mächtige Person zu sein, und mächtige Personen tendierten oft zu einem Hang für Willkür. Es schien ihm nicht falsch, auf die Symmetrie ihrer Situation hinzuweisen.

„Nun, Sheriff“, sprach Thorald, „Mir scheint das Ganze ebenso seltsam. Ich habe noch nie von einem Nottingham gehört. Wo befinde ich mich denn überhaupt?“

„Im Nottingham Castle, dem Herzen von England!“, kam die Frage wie aus der Pistole geschossen.

„England?“

Der Sheriff schüttelte lachend seinen Kopf: „Ich habe ja schon gehört, dass einen Schlag auf den Kopf einem Manne den Geist vernebeln kann, aber das ist schon unerhört, was ihr mir hier auftischt. Woher glaubt Ihr denn zu stammen, wenn ihr nicht einmal England kennt?“

„Wie ich bereits sagte: Ich stamme aus Andor, dem Drachenland.“

Der Sheriff verschluckte sich fast vor Glucksen, als er das hörte: „Drachen?! Als nächstes erzählt Ihr mir noch von Feen und Einhörnern!“

„Von den Drachen habt Ihr also schon gehört!“, rief Thorald freudig auf, „Dann werdet ihr ja wissen, dass die Drachenspezies vor langer Zeit schon ausgelöscht wurde bis auf den letzten. Dessen Hort... da direkt nördlich davon liegt Andor!“

Sheriff schlug sich auf die Schenkel: „Haltet ein, Thorald von Andor. Ich mag solche Geschichten so gerne wie der nächste, aber irgendwann wird es einfach zu viel. Ich hoffe, dass sich Euer Geist bald wieder beruhigt, auf dass wir einige Informationen über die Geächteten erfahren mögen.“

Wie konnte es sein, dass der Sheriff von Drachen, nicht aber von Andor gehört hatte? Hatte der Unterirdische Krieg vielleicht nur alle Drachen in der Nähe Andors ausgelöscht, aber weiter weg lebende verschont? Die Gerüchte über einen Eisdrachen im Hohen Norden würden dem ja auch entsprechen. Nicht, dass Thorald diesen Gerüchten Glauben schenken würde. Aber Ken mochte es, darüber zu diskutieren.

Da fiel Thorald etwas anderes auf: „Ihr behauptet, noch nie von Andor gehört zu haben. Wie kommt es dann, dass Ihr die andorische Sprache sprecht? Die Sprache, die den flüchtigen Ambacus aus Krahd einst von den Bewahrern gelernt wurde?“

Thorald besann sich, einst von einem legendären Trank gehört zu haben, welcher seinem Trinker erlaubte, Strukturen in gesprochener Sprache zu erkennen. War es vielleicht möglich, mit Runenmagie einen ähnlichen Effekt zu erzeugen? Runenmagie wie derjenigen, die den langen Gang zwischen Andor und dieser seltsamen Welt geziert hatte? Aber würde er die Sprache dieser Menschen dann so klar verstehen können? Und warum in diesem schrecklichen Akzent?

„Nein, mit ‚andorischer Sprache‘ hat das gar nichts zu tun. Ihr sprecht die Sprache Englands, wenn auch mit einer, wage ich zu behaupten, abscheulichen Betonung“, schmunzelte Sheriff.

War es möglich, dass sie tatsächlich dieselbe Sprache sprachen? Das würde, nein, musste bedeuten, dass sie gemeinsame Wurzeln hatten. Dass ein Kontakt zwischen den Welten bestand, oder dass gar Menschen aus Andor dieses Land besiedelt hatten, wenn nicht umgekehrt. Die Implikationen davon mochte sich Thorald gar nicht ausmalen. Nein, das hätten die Bewahrer vom Baum der Lieder mit all ihren Aufzeichnungen doch unmöglich unter Verschluss halten können.

Sheriff unterbrach Thoralds Überlegungen: „Auf welchem Pfad seid Ihr hierhergekommen? Zu Fuß? Zu Wagen?“

„Hoch zu Ross“, fuhr Thorald fort, „durch einen unterirdischen Gang. Erst hielt ich ihn für das Werk von Zwergen, aber die Zwerge hatten noch nie von einem England berichtet... vielleicht war das ein Feenpfad durch die Feenwelt? Solche können große Distanzen überwinden.“

„Ich sagte doch bereits: Haltet ein mit den Märchen!“, rief Sheriff nun mit einem bedrohlicheren Unterton in der Stimme, „Und bleibt bei der Wahrheit! Nichts von Feen, Zwergen, Magie oder dergleichen. Diese existieren nicht!“

Da erkannte Thorald endlich, dass Sheriff all diese Wesen und gar die Magie selbst für Hirngespinste zu halten schien, und verstummte überrascht. Sheriff war eindeutig ein reizbarer Mensch, und solange er Thoralds Erzählung nicht glaubte, war es vielleicht geschickter, sie nicht weiter auszuführen. Aber lügen wollte er auch nicht. Ach, wenn doch nur Ken hier wäre, der wüsste, was zu tun wäre.

Beim Gedanken an Ken spürte Thorald ein feines Stechen in seiner Brust, aber er ließ ihn nicht los, sondern versuchte, sich im Geiste Kens schnarrende Stimme vorzustellen:

„Ein sehr skeptischer und ignoranter Mensch ist das. Kein Wunder, dass er dir nicht glaubt, vor einem Jahrzehnt hättest du auch niemandem die Existenz von Tulgor abgekauft. Du musst etwas finden, was ihn überzeugt, dass du nicht auf den Kopf gefallen ist. Tische ihm keine alten heiteren Geschichten über Feen und Zwerge auf, sondern demonstriere ihm irgendetwas etwas Neues, Dunkles, vielleicht gar etwas Magisches?“

Da hatte Thorald eine neue Idee: Er griff in seine Hosentasche und zog den grünen Runenstein hervor, den er bei seiner Reise zu Elgas und Junas Hof gefunden hatte.

„Und was sagt Ihr zu diesem Prachtstück hier? Ein stärkender Stein, geprägt mit der Macht der Runen. Er war es, der mir den Gang in dieses Land überhaupt öffnete. Ein klassisches Beispiel für Magie!“

Sheriff griff nach dem Stein, wog ihn in der Hand und warf ihn dann zu Thorald zurück: „Ist das Futhorc? Ein hübscher Stein, kein Zweifel, aber ich erkenne darin keine Spuren von Magie.“

Jetzt, wo er es sagte, musste Thorald ihm recht geben. Der Stein glomm nicht mehr und als Thorald ihn an sein Ohr hielt, summte er auch nicht leise, wie die Runensteine es sonst taten. Er lag einfach nur still da. Als wäre die Magie in seinem Inneren erloschen.

Natürlich hatte Thorald schon erlebt, dass die Runensteine an verschiedenen Orten mit ihrer Umgebung mehr oder weniger resonierten – an den Eingängen zu den unterirdischen Zwergengängen waren sie beispielsweise mächtiger als über dem offenen Meer – aber dass er einen Runenstein fest in seiner Hand halten konnte, während dieser komplett erloschen blieb, das hatte er noch nie erlebt. Was war das hier für eine Umgebung?

„Keine weiteren Geschichtchen mehr?“, lachte Sheriff und behauptete stolz, „Das magischste, was ich je in meinem Leben erlebt habe, war eine präzise dressierte Krähe. Da draußen züchtet jemand Krähen und schickt sie auf uns los. Wenn wir eine von denen einfangen könnten, oder besser gar, zu ihrem Hüter verfolgen, dann wäre dies überaus hilfreich. Ich schwöre beim Herrn...“

„Beim Herrn?“

„Ach, kommt schon, Ihr glaubt doch nichts ernsthaft, dass ich Euch abkaufe, noch nie vom Herrn gehört zu haben.“

„Ich will Euch wirklich keinen Bären umbinden.“

„Bären hat es also dort, woher Ihr kommt?“

„Natürlich gibt es in Andor Bären!“

„Aber Gott kennt man dort nicht? Seid Ihr etwa ein Ketzer?“

„Ich habe noch nie in meinem Leben eine Kerze gezogen“, versicherte Thorald Sheriff.

„Ihr müsst Euch dringendst mal mit Father Egbert unterhalten, werter Thorald von Andor. Naja, sobald dieser sich wieder eingekriegt hat“, gluckste Sheriff, „Im Moment ist er nicht so gut auf die Krone zu sprechen, seitdem wir das goldene Kreuz der Kirche für die Kriegssteuer einziehen mussten.“

„Ach ja, Steuern. Das Volk mochte sie bei uns auch nicht sonderlich.“, murmelte Thorald bedrückt und erinnerte sich daran, wie Ken ihn zum Erheben größerer und größerer Steuern ermuntert hatte. Die Staatskasse Andors war nun praktisch am Überlaufen, aber das Volk war nicht besser auf die Krone zu sprechen. In den Flusslanden gab es gar Schreie nach Unabhängigkeit. Obwohl Thorald versprochen hatte, das Gold zum Wohle Andors zu nutzen, sobald sich eine Gelegenheit ergab.

Sheriff schein das nicht so zu sehen, denn er gluckste auf und stimmte zu: „Natürlich mag das Volk die Steuern nicht, die wollen ja lieber ihren Wohlstand für sich behalten. Aber Wohlstand der einzelnen Bürger besiegelt nun mal keine Bündnisse fürs ganze Volk. Und arme Schlucker sind abhängiger vom guten Willen der Krone. So können wir unsere schöne Macht beibehalten. Diejenigen, die nicht genügend leisten oder sich uns widersetzen, werden zu Geächteten, und diese wiederum erhalten früher oder später das, was sie verdienen: Den Tod!“

Thorald spürte ein flaues Gefühl in seinem Magen. Dieser Sheriff hatte nicht das Geringste Mitgefühl für seine Untergebenen. Daran erkannte man einen Tyrannen, das hatte ihm Brandur schon oft gesagt.

„Hört mal, werter Herr Sheriff...“, setzte Thorald an.

„Sheriff ist kein Name, sondern mein Titel“, brummelte der Sheriff, „Wenn ihr mich schon bei Namen nennen müsst, so nennt mich William von Wendenal.“

„Werter Herr von Wendenal“, setzte Thorald erneut an, doch blieb dann stumm. Sich mit dem Sheriff über die Moral seiner Taten zu unterhalten, wäre ein guter Weg, selbst auf die Feindesliste dieses Tyrannen zu landen. Stattdessen fuhr er fort:

„Ich könnte Euch vielleicht von der Wahrheit meiner Geschichte überzeugen, wenn ich Euch zum Höhlengang in mein Heimatland führte und dieses mit dem Runenstein öffnete. Wie klingt das in Euren Ohren?“

Sheriff lächelte mit kalten Augen: „Ha! Als ob ich so viel Zeit zur Verfügung hätte. Viele Angelegenheiten in der Grafschaft erfordern meine Aufmerksamkeit. Der hohe Prinz John hat sich höchstpersönlich hier in Nottingham einquartiert, könnt ihr es euch vorstellen? Was für eine Ehre! Seit der König in die Kreuzzüge gezogen ist, regiert der Prinz das Land an seiner Stelle. Darum müssen wir doppelt so hart durchgreifen, Wachpatrouillen verdoppeln und Räuber und Diebe ohne Gnade dem Tod überweisen. Und darum haben wir keinerlei Mannen für derlei Hirngespinste aufzuwenden.“

Thorald sagten eine Menge dieser Worte nichts, aber innerlich verfestigte sich sein Gefühl, dass er als König Andors nicht einmal so schlecht gewesen war. Zumindest hatte er nicht seine eigenen Untertanen ermorden lassen.

„So glaubt mir doch, warum würde ich Euch auf den Arm nehmen wollen?“, fuhr Thorald fort, der spürte, dass er ohne Beweise dieses Castle nicht mehr lebend verlassen würde, und der sich sicher war, dass ein Wiederholen seiner Behauptungen diese glaubhafter machen würde, „Genau dieser Runenstein hier öffnete den Ausgang zum Höhlengang zwischen unseren Reichen. Der war mit Runen verziert, Runen wie der auf diesem Stein hier. Der Ausgang hatte die Form eines riesiges steinernen Gesichts, und hinter seinem Mund liegt der Gang nach Andor...“

„Ein steinernes Gesicht?“ horchte der Sheriff plötzlich auf, „Hat es geweint?“

„Öh... vielleicht? Es war jedenfalls Wasser dort...“

Leise flüsterte der Sheriff: „‘\textit{Trittst du in den Kreis der Runen ein, wird der Strom der Tränen versieget sein.}‘ Ist das möglich, dass... ?“

Der Blick des Sheriffs, der Thorald nun traf, war urplötzlich kalt und entschlossen.

„Ich lasse einen Trupp bereitmachen. Führt sie zum steinernen Gesicht.“\bigskip







Thorald fühlte sich endlich wieder ein wenig königlich, als er auf einem edlen Pferd von Nottingham Castle aus aufbrach, in Begleitung zweier Wachen und des besten Söldners, den man für gutes Gold anheuern konnte. Guy von Gisbourne hieß er, und offenbar verdingte er sich als Jäger dieser Geächteten, die Thorald die mächtigen Schilde abgenommen hatten – wenn auch bislang bloß mit mäßigem Erfolg. Thorald zweifelte daran, dass ein so grimmiger Bursche dem geschickten Fenn aus dem Barbarenland im Fährtenlesen etwas vormachen konnte.

Der Sheriff betrachtete die Abreise des Trupps vom Torbogen aus, sein rotes Gewand in starkem Kontrast zum blauen Himmel. Insgeheim nahm Thorald sich vor, ein Wams in derselben Farbe in Auftrag zu geben. Die Farbe hatte etwas.

Auf dem Weg aus der Burg heraus blieb Thoralds Blick an einem großen Holzgestell hängen, an dem ein langes Seil mit einer Schlaufe am Ende hing. Thorald fragte Gisbourne, ob es sich dabei um ein Kunstwerk oder ein Werkzeug handelte. Gisbourne lehnte sich aus seinem Sattel und flüsterte Thorald ins Ohr: „Das ist der Galgen, an dem der Henker alle Gesetzesbrecher aufknüpft. Dieses Schicksal steht vielleicht auch dir bevor. Der Sheriff mag dir deine Wahnsinngeschichte abkaufen, aber so leicht trügt man mich nicht! Ich habe ein Näschen dafür, wo die Geächteten auftauchen, und diese Nase schlägt Alarm, wenn ich mich auch nur in deine Nähe komme. Sollte dein Tipp ins Leere schlagen, so lasse ich dich gefesselt und gekettet vor Prinz Johns Urteilspruch bringen. Und dann werde ich genüsslich zusehen, wie du am Galgen zappelst, bis dir die Luft ausgeht“

Thorald schluckte tief. Er hatte die letzten Jahre stark damit gekämpft, ein schlechter König zu sein, aber im Vergleich mit diesem Herrscher kam er sich ganz goldig vor. Ein Werkzeug zum Ermorden der eigenen Untertanen, so prominent zur Schau gestellt?! In was für einer Welt war Thorald hier nur gelandet?

„Meister Gisbourne...“, versuchte er, seine adlige Haltung zu bewahren, „Ich führe Euch hier freiwillig zu diesem Gang, in der Hoffnung, zurück in mein Land zu kommen. Ein wenig Dank wäre von Eurer Seite nicht fehl am Platze“

Thoralds zog rasch wieder den Kopf ein, als Gisbourne bloß knurrte und eine unwirsche Geste in seine Richtung machte. Was für eine grausame Kreatur dieser Söldner doch war. Thorald beharrte: „Ich kann Euch versichern, dass ich nicht lüge. Es mag sein, dass ich den Weg auf die Schnelle nicht finde, aber wenn wir nur lange genug suchen, werden wir die Stelle schon finden. Den Eingang kann man auf jeden Fall nicht übersehen. Seht, wir folgen den Radspuren des Wagens, der mich zur Burg gebracht hat. Folglich liegt der Gang in dieser Richtung. Haltet einfach Ausschau nach einem Runenkreis oder einem steinernen Gesicht.“

Thorald hob seine Hand und zeigte grob in die Richtung, in welcher er das Gesicht vermutete. Ein Rascheln ließ seinen Blick ins Unterholz gleiten, wo sich soeben eine grau gewandte Gestalt mit einem großen Buckel in den Schatten zurückzog. War das Reka, die Kräuterhexe? War sie etwa auch dem Tunnel in dieses fremde Land gefolgt?

„Keine Ablenkungen“, zischte Gisbourne, „Auf, auf! Die Geächteten kommen näher, das spüre ich. Wenn wir uns beeilen, können wir sie vielleicht überraschen.“

Da fühlte Thorald, als hätten die Baumwipfel Augen und würden auf den kleinen Trupp des Königs starren, wie er durch den Sherwood Forest ritt. Ein Blick nach hinten verriet ihm, dass die beiden Wachen ebenfalls vorsichtig die Bäume inspizierten. Nur Guy von Gisbourne ritt mit starrem Blick nach vorne voran.

Sie waren kaum eine Viertelstunde weiter gezogen, da schrie Gisbourne plötzlich auf. Thorald hatte seine Schwierigkeiten damit, sein Pferd zu zügeln, da zischten auch schon mehrere irgendetwas an seinem Kopf vorbei.

Im Licht, das die sinkende Sonne durch die Wipfel warf, erkannte Thorald mit Schrecken, wie die beiden Wachen hinter ihm zu Boden sanken. Aus der einen ragten zwei Pfeile, von denen der eine den anderen beinahe vollständig gespalten hatte.

Thorald wurde bleich. Hilflos suchte er die Umgebung ab, nach einem Zeichen, wer denn hinter diesem Überfall steckte. Gisbourne starrte die meisterhaft geschossenen Pfeile an und flüsterte: „Das ist Robin Hood. Das ist ihr Anführer!“

Ein irres Glänzen schlich sich in seine Miene, als er ausrief: „Zeig dich, Robin, wo versteckst du dich?“

Als wäre es seine Antwort, surrten zwei weitere Pfeile durch die Luft und schlugen gegen die Brustplatte des Söldners. Gisbourne keuchte, aber die Pfeile hatten seinen Brustpanzer nicht durchschlagen und fielen wirkungslos zu Boden.

Gisbourne zog einen verdreckten Helm hervor und setzte ihn auf. Lilane Federn ragten stolz daraus hervor und nur noch seine wahnsinnig glänzenden Augen waren von ihm zu erkennen. Jetzt würde es erst recht schwer werden, ihn zu treffen.

Thorald fluchte und glitt von seinem Pferd, in der Hoffnung, am Boden vor einem Pfeilhagel sicherer zu sein. Er landete unsanft in einer Ansammlung von Waldpilzen, welche prompt zu stäuben begannen und seine Nase reizten. Thorald nießte einmal, zweimal laut und rollte sich aus den Pilzen auf den matschigen Waldboden. Es stank nach Pferdedreck.

Als Thorald sich aufrichtete, sah er zwei Gestalten entgegen, welche in die hell erleuchtete Lichtung traten. Die eine war schlank, aber großgewachsen, in einen himmelblauen Umgang samt Kapuze gekleidet, und hielt einen mächtigen Bogen auf Gisbourne gespannt. War das dieser legendäre Robin Hood, der Anführer der Geächteten? Die andere Gestalt trug ihr langes feuerrotes Haar offen und stützte sich auf eine schwere Axt. Wenn Thorald sich nicht täuschte, wehte der Wind gerade ein bisschen zu viele Blätter um sie herum, als dass es ein Zufall sein könnte. Ein schwarzer Rabe stürzte sich aus dem Himmel und flog der Gestalt ohne Kapuze in ihre ausgestreckte Hand, woraufhin sich Thorald an die Worte des Sheriffs erinnerte:

„Da draußen züchtet jemand Krähen und schickt sie auf uns los.“

„Zurück!“, erklang ihre helle, gebieterische Stimme, „Ziehet von hinnen und wir schenken Euch Euer Leben!“

„Zur Hölle! Wo ist Robin Hood?!“, zischte Gisbourne und presste seine Fersen in die Flanken seines Pferdes. Dieses wieherte, richtete sich auf seine Hinterbeine und sah ganz gefährlich aus. Dann ließ es seine Vorderbeine wieder auf den Boden prallen und nahm geifernd Kurs auf die beiden Geächteten. Gisbourne zückte seine blanke Klinge und richtete sie auf die Gestalt mit dem Raben.

Der himmelblaue Bogenschütze feuerte auf den Söldner, einmal, zweimal, aber wieder prallten die Pfeile nutzlos an Gisbournes Rüstung ab. Mit einem frustrierten Grunzen schulterte der Bogenschütze seinen Bogen wieder und rief: „Gisi? Darf ich dich Gisi nennen? Die Welt wird sich erfreuen an den Spottliedern, die ich über dich verfassen werde... sobald du erst einmal das Zeitliche gesegnet hast!“

Der Bogenschütze zog ein langes Messer hervor und richtete es auf den anstürmenden Gisbourne, scheinbar ungerührt von dessen Ansturm. Gisbourne gab sich ähnlich ungerührt, zügelte sein Pferd und bellte vom Rücken seines Gauls einen knappen Befehl. Wie aus dem Nichts tauchten hinter den beiden Geächteten zwei weitere rote Wachen auf und fuchtelten mit langen Speeren. Die Frau mit dem Raben wich zurück und sprang auf den nächsten Baum, während ihr Rabe der Wache ins Gesicht flog und allen Anschein nach ansehnlichen Ärger bereitete.

Dann aber wurde der Rabe davongescheucht und die beiden Wachen drängten den übriggebliebenen Bogenschützen in Thoralds Richtung. Ein Speer verfing sich in der himmelblauen Kapuze, riss sie herunter und enthüllte ein faltenloses Gesicht mit einem strubbligen schwarzen Haarschopf. Das war doch noch kaum ein erwachsener Bursche! Thorald sah zu, wie der Geächtete gefährlich unkontrolliert mit seinem langen Messer herumfuchtelte, und beschloss, einzugreifen. So griff er nach einem am Boden liegenden Büschel von Ästen und schleuderte diese ungelenk von sich. Nicht viele der Äste trafen ihr Ziel, doch zu seiner Freude sorgte sein Eingreifen dafür, dass der Bursche sich umdrehte und die beiden Wachen sich auf ihn stürzen konnten. Nach einem kurzen Gerangel wurde das Messer des Burschen fortgeschleudert. Gisbourne sah dem ganzen Vorgehen aus einiger Entfernung grinsend zu. Thorald richtete sich zu seiner vollen Größe auf und wischte sich den Dreck von seinen Kleidern.

Ein dumpfer Schlag direkt neben ihm ließ sein Herz wieder in die Hose sinken.

„Haste mich vermisst?“, ertönte eine tiefe Stimme. Der Hüne, der Thorald erst vor Kurzem so freundlich in diesem Reich begrüßt und ihm die beiden mächtigen Schilde entwendet hatte, war auch hier! Thorald schrie mutig auf und gab Fersengeld, aber kaum zwei Mannslängen weiter erwischte ihn der kräftige Griff des Hünen dennoch.

Der Hüne warf Thorald ohne Mühe über seine Schulter. Jetzt konnte Thorald nur noch den verschwommenen Waldboden ausmachen, wie er hin- und herschaukelte, während der Hüne mit großen Schritten davonrannte. Ein mächtiger Satz, dann ein Knarzen wie von einem Baumstamm, dann... Stille. Ein Schrei erklang, in dem Thorald die Stimme Gisbournes zu erkennen glaubte, aber er schien weit weg.\bigskip







Thorald wurde sanft zu Boden gelassen. Das war kein Waldboden, sondern... Bretter? Eine Brücke in den Bäumen? Ehe Thorald sich genauere Gedanken dazu machen konnte, hatte sich der Hüne bereits über ihn gebeugt und platzierte seine breite Hand fest auf Thoralds Mund.

„Hübsch still sein, gell?“, wies er Thorald an. Dieser hustete und blickte sich panisch um, aber Widerstand schien ihm im Moment nicht das geschickteste Vorgehen zu sein.

Es raschelte neben ihm und die Gestalt mit dem Raben trat in sein Blickfeld. Überrascht erkannte Thorald, dass es sich um eine Dame in einem langen Kleid handelte. Das Bild passte so gar nicht zur mächtigen Axt, die sie mit zwei Händen führte.

„Sie haben Will geschnappt“, fluchte sie.

„Un‘ wir ha’m den hier“, meinte der Hüne,

„Aber wir hängen unsere Gefangenen im Gegensatz zu ihnen nicht einfach so, wenn sie nichts dagegen tun.“

Der Hüne kratzte sich am Kinn und antwortete dann: „Wenn Gisi ihn geschnappt hat, gibt es so hurtig nichts, was wir für ihn tun könn‘. Komm, Marian, wir bringen den hier zu Robin, und dann guck’n wir weiter, wie wir Will aus dem Castle kriegen.“

„Deine Hoffnung hätte ich gerne“, schüttelte die Dame – Marian – ihren Kopf.

Der Hüne meinte schlicht: „Der Herr ist mit uns.“

Marian schluckte tief, nickte dann aber: „Na gut, dann auf zu Robin mit ihm. Zuletzt hat er auf der anderen Seite der Lichtung Position bezogen. Suchen wir ihn dort.“

„Nicht nötig“, erklang eine leise Stimme dicht neben dem Trio, und alle drei zuckten zusammen, wenn auch Thorald am stärksten. Er versuchte erfolglos, seinen Kopf so zu drehen, dass er den Neuankömmling erkennen konnte.

Der Neuankömmling indes beugte sich so nahe an Thorald heran, dass dieser seinen Atem spüren konnte, und murmelte: „Ich bin Robin von Loksley.“

Thorald stammelte ein „Sehr erfreut“, das durch die Hand des Hünen kaum zu verstehen war.

Robin fuhr ungerührt fort: „Wir sind etwas in Eile, darum komme ich gleich zum Punkt: Wer seid Ihr, und warum wisst Ihr, wo unser geheimes Lager liegt?“

Thorald schluckte. Er wusste doch gar nicht, wo ihr geheimes Lager lag. Und so skeptisch, wie der Sheriff auf seine Geschichte reagiert hatte, war es vielleicht nicht das Geschickteste, einem Haufen Halunken seine Herkunft zu erklären versuchen.

Schon entfernte der Hüne seine Hand von Thoralds Mund. Thorald krabbelte ein wenig zurück und richtete sich auf. Ja, er stand auf einer Art Baumbrücke, durch deren Ritzen er den erstaunlich weit von ihm entfernten Waldboden erahnen konnte. Ihm wurde etwas schwindelig und er richtete seinen Blick lieber auf die drei Geächteten, die ihn gespannt anblickten. Neben ihnen sah Thorald nun den Sack, den der Hüne Thorald bei ihrer ersten Begegnung abgenommen hatte. Die beiden mächtigen Schilde waren hier, direkt in seiner Reichweite!

„Nun?“, fragte Robin erneut. Sein Ton war nicht freundlicher geworden.

„Ich... ich bin ein Reisender aus einem fremden Land“, gab Thorald vorsichtig von sich, „und ich habe wirklich nicht die geringste Ahnung, wo Euer geheimes Lager liegt. Und ebenso wenig Interesse daran.“

Der Hüne gluckste: „Wir würd’n dir ja gerne glauben, aber dass du den groben Gisi direkt in uns’re Richtung geführt has‘, sagt ‘was andres als deine Stimm‘.“

Robin nickte: „Wenn der Hüter des Waldes mich nicht vorgewarnt hätte, hättest du Guy und seine Truppen direkt zu uns gebracht.“

„Aber... aber ich bin doch erst seit einigen Stunden hier! Wie soll ich denn ein geheimes Lager gefunden haben, welches nicht einmal der beste Söldner des Landes aufspüren konnte?“

Robin wandte sich den anderen beiden zu: „Er lügt nicht... könnte es sein, dass er vielleicht ebenfalls...“

Marian unterbrach ihn: „Was ist das für ein fremdes Land, aus dem du stammst?“

Ach, was soll’s. Etwas ausdenken wollte Thorald sich nun wahrlich nicht, das war nie seine Stärke gewesen.

So begann er: „Mein Heimatsland heißt Andor. Ihr mögt noch nicht von ihm gehört haben, aber es existiert. Wirklich. Und ich bin dort der K... eine wichtige Persönlichkeit. Also wagt es nicht, mir etwas anzutun. Ihr wollt keinen Ärger mit uns anfangen.“

Seine Drohung ignorierend, sahen sich Robin und Marian wissend an, und der Hüne gluckste erneut. Dann sprach er: „Witzig, dass du‘s so sagst. Wir ha’m den Namen nämlich g’rad vorhin schon mal g’hört.“

Als wäre das ihr Stichwort gewesen, schwang sich eine dunkel gewandete Gestalt auf die Plattform in den Bäumen und kam neben Robin zum Stehen.

Das Licht der untergehenden Sonne traf auf einen kahlen Schädel. Thorald erkannte überrascht Ken Dorr, seinen treuen Ratgeber, wie er ungerührt neben den drei Geächteten stand und ihnen Rapport gab: „Sie haben Will ins Castle gebracht. Ich habe es mit eigenen Augen gesehen.“

Robin erwiderte ebenso ungerührt: „Dann hast du bestimmt nichts dagegen, wenn ein paar fremde Augen das bestätigen lassen.“

Er gab ein Handzeichen. Marian nickte und sprang von der Baumbrücke. Eine Rabe stürzte aus dem Himmel und folgte ihr.

Ken Dorr fixierte Robin aus kalten Augen und lachte auf: „Endlich jemand mit Verstand, Robin von Loksley! Es wäre töricht gewesen, meinem Rat einfach so zu vertrauen. Schön, zu sehen, dass...“

„Still, Kennard. Zu dir komme ich noch.“

Ken verstummte und zog sich ein bisschen zurück, die Hand scheinbar locker auf seiner Hüfte liegend, aber gefährlich nahe an seiner Schwertscheide. Sein Gesicht war wie eine grimmige Maske. Er gab nicht zu erkennen, ob er Thorald erkannte. Aber es musste Ken Dorr sein, diese Narbe an seiner rechten Wange würde Thorald im Schlaf wiedererkennen.

„Das is‘ gut“, meinte der Hüne nun, „Solang‘ Will nicht in dieser verfluchten Festung Blackgarden steckt, ha’m wir ‘ne gute Chance, ihn da wieder ‘rauszuhauen.“

Robin nickte und tigerte auf der kleinen Brücke hin und her, während er sich am Kopf kratzte:

„Wenn Kennards Bericht stimmt, müssen wir uns beeilen. Irgendwie in die Burg schleichen, bevor Prinz John zu Will kommt – oder der Henker. Wir brauchen ein Seil.“

„Colin der Zimmermann könnte noch ein‘s bei sich haben.“

„Wollen wir die Dorfbewohner wirklich noch mehr in Gefahr bringen?“

„‘S‘is immerhin Will, um den’s hier geht.“

Thorald blickte zu Robin und dem Hünen, wie sie verschiedene Pläne besprachen. Wie gemeine Diebe wirkten sie nicht, erst recht nicht, als sie die Sicherheit der Dorfbewohner zu besprechen begannen und festlegten, wie viel Proviant sie den Hungernden verteilen konnten. Thorald schlich sich langsam rüber zu Ken. Dieser wirkte im Gegensatz zu den Geächteten wie ein waschechter Gauner, wie er mit einem schiefen Lächeln auf den Lippen mit einer Goldmünze umherspielte. Thorald lag immer noch ungut im Magen, was in Andor zwischen ihm und Ken vorgefallen war. Aber darum konnte er sich auch später noch kümmern. Im Moment war es einfach nur eine unglaubliche Erleichterung, ein bekanntes Gesicht zu sehen.

Ken versuchte eine Zeit lang, Thorald zu ignorieren, aber als dieser zu nahe an ihn getreten war, setzte er ein gequältes Lächeln auf uns sprach theatralisch:

„Kennard von Dorr ist mein Name. Fremder in einem fremden Land, wie du, so scheint es.“

Endlich fiel der Groschen bei Thorald, dass Ken nicht preisgeben wollte, dass sie sich bereits kannten, und er nickte theatralisch:

„Thorald ist der meine. Ummm... Ihr... Ihr stammt ebenfalls aus Andor?“

Ken verdrehte die Augen ob dem fehlenden Schauspieltalent, nickte aber brav.

„Ja, zu euch beiden und euren seltsamen Namen kommen wir noch“, wandte sich Robin ihnen beiden wieder zu, „Aber Will aus der Gefangenschaft zu befreien hat im Moment Priorität.“

Der Hüne nickte zustimmend.

Ein Rabe kam angeflogen und krähte zweimal.

„Nun denn, sieht so aus, als ob du die Wahrheit gesagt hast, Kennard. Will wurde ins Nottingham Castle gebracht“, nickte Robin, „John und ich machen uns auf den Weg dorthin. Kennard, kannst du dich um unseren Adligen kümmern, bis wir zurückkommen? Danach entscheiden wir, was wir mit ihm und seinem Gespür für versteckte Lager tun wollen.“

Ken schien kurz etwas erwidern zu wollen, nickte dann aber brav und legte Thorald possessiv eine Hand auf die Schulter. Thorald schauderte leicht.

Robin packte den Sack mit den mächtigen Schilden. Dann sprangen er und John von der Baumbrücke – es ertönten ein leiser und ein mächtiger Aufprall – und zogen davon. Ken wartete einige Augenblicke und wandte sich dann mit einem breiten Grinsen Thorald zu:

„Ich nehme alles zurück, was ich über seine Weisheit sagte. Es war töricht von ihm, uns zwei Fremde zusammen zurückzulassen.“

Es war, als würde eine große Last von Thoralds Schultern fallen. Das hier war Ken, wie er ihn kannte und liebte. Der weise Ken, der stets einen Plan hatte und die Lage zum Bessern wenden konnte.

„Nun, du hast ja auch so getan, als würden wir uns nicht kennen. Und sie haben gerade Wichtigeres im Kopf“, relativierte Thorald. Resigniert ergänzte er: „Wegen mir ist dieser Will in Gefangenschaft geraten und wird jetzt vor den Folterknecht und dann den Henker kommen. Ich habe wieder einmal alles falsch gemacht.“

„Machst du Witze?“, rief Ken, „Das war doch ihre Schuld. Sie dachten, dass du sie direkt zu ihrem Lager führst, obwohl du dieses noch nicht einmal gesehen hast. Allein aufgrund der Warnung eines uralten Waldverrückten! Wer aufgrund einer derart geringen Sachlage bereits zu einer Tat schreitet, ist ganz und gar selbst für deren Scheitern verantwortlich.“

Thorald war die Spitze in Kens Worten nicht entgangen, die sich nur allzu gut auf Reka übertragen ließ. Rekas Warnung bezüglich Ken fiel ihm wieder ein. Er ignorierte sie geflissentlich, wie es einem König gebührte, unterdrückte seine Schuldgefühle, wie er es schon so oft getan hatte, und stellte stattdessen eine ebenso dringende Frage: „Wie... wie hast du mich überhaupt gefunden?“

Ken gluckste: „Die Hufspuren im Schlamm vor der Quelle des Likko waren kaum zu übersehen. Du bist nicht wirklich geübt im Verstecken, oder? In der Höhle musste ich dann nur der Spur des tropfenden Pferds folgen und bin durch einen langen Gang zu diesem steinernen Gesicht gekommen. Hat eine Zeit lang gebraucht, bis ich den Ausgang öffnen konnte. Zum Glück hatte ich einen Casamatuc zur Hand.“

Ken Dorr schüttelte ein kleines Werkzeug aus seinem Ärmel und überreichte dieses Thorald zur Inspektion. Tatsächlich, das könnte ein Casamatuc sein! Das waren Zwergenwerkzeuge, mit denen man so gut wie jedes Schloss öffnen konnte. Sie waren relativ selten, Thorald hatte bislang nur in Erzählungen von ihnen gehört. Wie Ken an diesen Casamatuc gekommen war, wollte Thorald lieber nicht nachfragen.

Während Ken ihm fröhlich den Casamatuc zeigte, konnte Thorald nicht umhin, sich zu wundern, wie Ken so gelassen sein konnte. Der zornerfüllte Schemen, der Ken im Eingang zu Elgas Hütte gewesen war, schien vollkommen verschwunden. Oder zumindest unterdrückt. Thorald wollte lieber nicht zu lange darüber nachdenken, sondern führte stattdessen seine eigene Geschichte aus.

„Ich hatte natürlich keinen Allzweckschlüssel, doch bemerkte ich tatsächlich erst, dass das steinerne Gesicht sich verschließen kann, nachdem ich hindurchgetreten war. Der Mund scheint sich automatisch zu öffnen, sobald man einen Runenstein in seine Nähe bringt. Der ganze Gang war ja auch mithilfe von Runenmagie beleuchtet. Aber kaum war ich ein bisschen weitergezogen, wurde ich auch gleich von diesem Hünen, John, überfallen und dieser hat mir... mit... mit einem fiesen Trick die Schilde abgeknöpft. Wie... wie hast du dich bei den Geächteten eingeschlichen?“

Kens Blick trübte sich.

„Ach Thorald, ich habe dir doch schon einmal davon erzählt. Ich musste mich in meiner Kindheit in der Gasse rumschlagen, und habe da das eine oder andere Talent aufgegriffen. Den gefürchigen Dieb zu spielen, gehörte leider dazu. Dieser Will Scarlet ist wie ein Waldgeist vor mir aufgetaucht und wollte mich überfallen, aber da habe ich ihm gezeigt, dass man Ken Dorr nicht so einfach beraubt.“

Ken grinste kurz, dann wurde seine Miene aber wieder finster.

„Ein Glück, dass ich dich so schnell finden konnte. Ich werde nicht lügen, Thorald, dein irrsinniger Umgang mit den mächtigen Schilden hat mich ziemlich wütend gemacht. Aber jetzt müssen wir uns gerade auf Wichtigeres fokussieren. Dieses Steingesicht, durch welches wir hierher kamen, liegt nicht weit von hier. Wir können zurück nach Andor und mit der gesamten Rietgarde zurückkehren, um uns die Schilde zu erkämpfen.“

Der Gedanke war verlockend. Thorald vermisste das Rietland schon jetzt. Dort könnte er sich endlich wieder sicher fühlen. Saubere Kleider anziehen. Sich weder mit tyrannischen Herrschern noch Geächteten herumschlagen. Nur... nur mit Ken. Aber dann müsste Thorald auch mit leeren Händen zurückkehren, während die mächtigen Schilde in diesem unbekannten Land saßen und weiß Mutter Natur was anstellten. Und die Rietgarde von der Rietburg abzuziehen war ein Fehler, den Thorald nur einmal machen würde. Die erste Befreiung der Rietburg hatte genug Leben gekostet. Ganz zu schweigen davon, dass eine fremde Armee im Sherwood Forest den Bewohnern von Nottingham Castle sicherlich nicht gefallen würde, und ihre Wachen waren nicht schlecht ausgerüstet, nach dem wenigen, was Thorald im Castle hatte erkennen können. So viel Verstand von der Kriegskunst hatten ihm seine Lehrmeister eingebläut.

Dann schoss Thorald wieder das Bild durch den Kopf, wie dieser junge Geächtete, Will, wegen Thoralds Angriff strauchelte und von den roten Wachen zum grimmig grinsenden Guy von Gisbourne geschleift wurde. Nein, Thorald würde nicht fliehen. Ganz abgesehen von den mächtigen Schilde hatte er durch seine Taten Will dem Tode verurteilt. Dem unrechten Tode. Als König war es seine Pflicht, Unrecht entgegenzuwirken!

Thorald schüttelte seinen Kopf: „Ken, das Volk hier leidet unter einem tyrannischen Herrscher. Ohne Will fehlt den Geächteten ein wichtiger Streiter. Will wird wegen mir vor eine Mordmaschine kommen. So etwas habe ich meiner Lebtag nicht gesehen. Und im Gegensatz zu den Geächteten habe ich einen sicheren Zugang in die Burg. Nur ich kann ihm jetzt noch helfen.“

„Ist bei dir noch alles in Ordnung im Oberstübchen?“, spottete Ken, „Du willst mir doch nicht ernsthaft sagen, dass du dich jetzt um eine dahergelaufene Bande von Aufmüpfigen scherst! Du suchst doch nur wieder nach einer Möglichkeit, deinen Stolz zu bewahren! ‚Thorald, der gute und gerechte König‘. Diese Kutsche ist doch schon lange abgefahren. Hör auf, dir etwas vorzumachen, und komm mit mir zurück nach Andor, lass die Rietgarde dies handhaben!“

Thorald blieb trotzig: „Meine Meinung steht fest. Wir sind... ich bin Schuld daran, dass Will gefangen genommen wurde. Ein junger Mann ist das, ein guter Mann, der noch sein ganzes Leben vor sich gehabt hatte. Es ist meine Pflicht, ihm da wieder rauszuhelfen.“

Ken griff sich bloß an den Kopf und stöhnte.

„Na komm“, meinte Thorald, und klopfte Ken auf die Schulter, „Wir haben uns lange genug vom Schicksal herumschubsen lassen. Lass uns aktiv werden.“








\newpage
\section{Thorald tut etwas Mutiges}

Die Nacht war bereits lange angebrochen und ihre finsteren Schatten schmiegten sich an die Mauern von Nottingham Castle, als Thorald den Eingang erreichte.

„Man... man öffne die Tore! Ich bin entkommen!“, entsprangen die krächzenden Worte seinen aufgerissenen Lippen.

Eine Wache in ihrem typischen roten Gewand beugte sich über den Torbogen und verschwand dann gleich wieder dahinter. Thorald hört leises Getuschel von oben und versuchte, noch ein wenig mitleiderregender auszusehen.

„Ist das etwa dieser elende Bettler?“

„Bist du blind? Der ist doch viel älter.“

„Wir können doch nicht einem dahergelaufenen...“

Thorald erkannte, dass er seine Taktik ändern musste. Er richtete sich auf, so gut es ging, warf sich in Pose und ließ mit tönender Stimme verlauten:

„Mein Name ist Thorald vom Königreich Andor, und ich habe Informationen über den Aufenthaltsort der Geächteten! So lasset mich ein!“

Wieder beugten sich neugierige Wachen-Gesichter über den Torbogen, diesmal mehrere.

Schließlich hörte Thorald eine Stimme: „Holt den Sheriff!“

Er entspannte sich. So weit, so gut.\bigskip







Der Sheriff war ganz und gar nicht glücklich darüber, mitten in der Nacht geweckt worden zu sein. Er trug ein vollkommen rotes Nachthemd mit einer eleganten Mütze (Thorald merkte sich in Gedanken die Farbe, damit er in Andor selbst ein solches Gewand in Auftrag geben konnte) und rümpfte die Nase, als der mit Schlamm bedeckte Thorald in sein Arbeitszimmer spazierte. Seine Augen leuchteten aber rasch auf, als Thorald von einem fantastischen erfundenen Gefecht mit den Geächteten berichtete, an dessen Ende Thorald knapp mit seinem Leben davonkam – und nun eine noch genauere Vorstellung davon hatte, wo sich das geheime Lager der Geächteten befinden könnte. Gleich morgen in aller Herrgottsfrühe würde Guy von Gisbourne einen neuen Suchtrupp anführen, versprach der Sheriff.

Er packte zwei Karaffen Wein aus einer Truhe hervor und bot eine davon Thorald an.

Dieser lächelte: „Jetzt sprecht ihr definitiv meine Sprache!“\bigskip







Der Heiler von Nottingham Castle war auch nicht sonderlich glücklich darüber, aus dem Schlaf gerissen zu werden, aber wie könnte er dem Sheriff einen Wunsch abschlagen? Thorald wurde gründlich auf Wunden untersucht und schließlich durchgewinkt. Daraufhin ließ der Sheriff ihm ein Bad bereiten und ein frisches Kleidungsset auslegen. Endlich, endlich, konnte Thorald sich etwas entspannen und die Ereignisse des letzten Tages Revue passieren lassen. So kam es, dass er fast einnickte, ehe er sich an den wahren Grund erinnerte, weswegen er sich zurück ins Nottingham Castle geschlichen hatte. Dieser Grund saß wahrscheinlich gerade einige Stockwerke über ihm in einer klammen Zelle und bibberte dem Tod entgegen. Nun, hoffentlich nicht mehr lange.

Thorald bedankte sich ausschweifend beim Sheriff von Nottingham für das Bad und die Kleidung und fragte dann genauso schwungvoll nach einer Schlafgelegenheit. Der Sheriff gähnte todmüde und wies gedankenverloren auf einen kleinen Häuserkomplex: „Da drin wird man bestimmt angemessene Unterkunft für dich finden können.“

Thorald bezweifelte das, aber das war im Moment auch nicht wichtig. Er bewegte sich auf das Haus zu und ließ seinen Blick sorgfältig über die Burgmauern wandern. Die wenigen Wachen, die dort Wache hielten, beachteten das Innere der Burg nicht. Der Sheriff zog sich in seine Privatkammer zurück. Der Weg war frei.

Nottingham Castle war etwa von derselben Größe wie die Rietburg, aber die Architektur dieser Burg erschien ihm dennoch verwirrend. Thorald brauchte einige Versuche, um eine offene Tür in den großen Turm zu finden, in dem Ken den Gefangenentrakt vermutet hatte.

Schlimmstenfalls, falls er entdeckt wurde, könnte er sich auf Schlafwandeln berufen. Oder darauf, dass er nur den Abort suchte?

Glücklicherweise musste Thorald keine halbherzigen Ausreden einsetzen, denn keine Menschenseele kreuzte seinen Weg, bis er endlich vor einer schweren Eichentür stehen blieb, welche vielsagend mit einer massiven Eisenkette verschlossen war. Dahinter hörte Thorald ein leises Gemurmel. Er versuchte, einen Blick durch die Spalten zu werfen, aber es fiel ihm denkbar schwer.

Schlussendlich wagte er es, seine Stimme erklingen zu lassen, lehnte sich gegen die Tür und wisperte: „Wer ist da?“

Stille. Überraschte Stille? Thorald überlegte bereits, weiterzuziehen, als eine tiefe Antwort aus dem Innern erklang: „Hier sitzt Will Scarlet.“

Innerlich triumphierte Thorald. Äußerlich fischte er mit zitternden Fingern den Casamatuc aus seiner Tasche, das Zwergenwerkzeug, mit dem Ken durch das steinerne Gesicht in diese Welt hatte treten können.

„Halte still, Will, ich hole dich hier raus“, flüsterte Thorald durch die Zellentür. Ein Klick, zwei Klicks, beim dritten klemmte es, etwas Rütteln, ein Klack, eine Drehung, und schon hatte der Casamatuc sein Werk vollbracht und das schwere Eisenschloss fiel polternd zu Boden. Thorald vernahm ein Ächzen von der anderen Seite der Tür und dann ein heller Aufschrei: „Nein, nicht!“

Diese Stimme war viel höher gewesen als die tiefe Stimme, die ihm vorhin geantwortet hatte. Thorald musste sich nicht lange Gedanken über dieses Rätsel machen, denn schon wurde die Tür vor ihm aufgerissen und der König von Andor blickte dem Urheber der tiefen Stimme in die glühenden Augen. Das musste er sein: Prinz John, von dem der Sheriff gesprochen hatte. Der neue Herr von Nottingham Castle, mit einem blutigen Messer in der Hand, welches unschön gut zu zwei Schnitten im Gesicht des jungen Mannes passte, der sich am anderen Ende der Zelle ins Stroh kauerte.

„Nein, rette dich!“, rief der junge Mann mit der schwarzen Strubbelfrisur nun erneut. Er trug immer noch einen himmelblauen Umhang. Das war unzweifelhaft Will Scarlet, der Bogenschütze, für dessen Festnahme Thorald gesorgt hatte.

„Nun, das ist ja eine hübsche Überraschung“, grinste Prinz John süffisant, „Und... sind das etwa Kleider des Sheriffs? Ihr seid mir so manche Erklärung schuldig. Aber keine Sorge, die werdet Ihr mir liefern.“

Wie Prinz John da auf ihn zuschritt, mühelos mit dem kleinen Foltermesser in der Hand spielend, fühlte sich Thorald plötzlich an die Zeit als kleiner Junge erinnert, als er zum ersten Mal einem Troll gegenübergestanden war. Und ihm wurde schmerzlich bewusst, dass er wieder einmal keine Waffe bei sich trug.

Aber das hier war anders. Thorald war kein kleiner Junge mehr. Sein Gegenüber war kein Troll, sondern bloß ein bösartiger Mensch. Und Thorald brauchte keine Waffen, um sich zu verteidigen. Diesmal würde er nicht wegrennen.

So nahm Thorald all seinen Mut zusammen, hob seine Fäuste und schlug nach der Magengegend des Prinzen. Es war kein schöner Stil, und Harthalt wäre nicht stolz darauf gewesen, wie er seine Deckung vernachlässigte, aber es reichte, um den überraschten Prinzen keuchend zurückzudrängen. So waren manche Prinzen halt. Viel Worte, weniger dahinter.

„Renn, Will, renn!“, rief Thorald.

Will rappelte sich auf und grinste schwach: „Ich habe die letzten Stunden in dieser engen Zelle nicht umsonst Kräfte gespart. Lass dich nicht abhängen!“

Schon war Will an Thorald vorbei und aus der Tür geflitzt. Gute Güte, mit dieser Geschwindigkeit könnte der Bursche einem Skral einen ansehnlichen Wettkampf bieten!

Thorald war so beeindruckt, dass er beinahe selbst vergaß, die Beine in die Hand zu nehmen. Beinahe. Der hasserfüllte Blick Prinz Johns war am Ende doch genug, ihn aus seiner Starre zu reißen und Reißaus nehmen zu lassen.

Weiter vorne am Burgtor regten sich die Wachen, als sie erkannten, dass im Innern der Burg Spannenderes lief als außerhalb – erst recht, als Prinz John aus der Gefangenzelle zu taumeln kam und Zeter und Mordio schrie.

Will drehte sich abrupt um und rannte zurück, auf Thorald zu: „Gibt es einen Plan? Ist Robin hier?“

Thorald versuchte, wieder zu Atem zu kommen, und winkte ab. Wie lange war es nun schon her, dass er an einem Reitgefecht teilgenommen hatte? Wann war er so außer Kondition geraten?

Will ließ die Schultern sinken: „Es sieht nicht rosig aus. So einfach kommen wir nicht von der Burgmauer runter. Falls ich das überlebe, werde ich ein Lied auf deinen törichten Todesmut machen.“

„Törichter Todesmut ist richtig“, bekräftigte Thorald, „Ich lasse nicht zu, dass du meinetwegen umkommst.“

Will schaute ihn verdutzt an und schien sich unsicher, was er mit dieser Aussage anfangen sollte. Ehe er zu einer Entgegnung ansetzen konnte, packte Thorald ihn kurzerhand und warf sich gemeinsam mit ihm über die Brüstung der Burgmauer.

Der Schreck stand Will immer noch ins Gesicht geschrieben, als er sich hustend und prustend aus dem Heuwagen an der Mauer schälte. Ein Glück, dass Ken wie versprochen den Heuwagen hierhin schieben hatte können.

Will setzte an, das Weite zu suchen, doch Thorald hielt ihn zurück.

„Warte!“

„Was?! Wenn wir hier bleiben, dann werden wir gefunden.“

„Sie werden uns nicht sehen... hoffentlich. Sie werden annehmen, dass wir das Weite gesucht haben.“

„Aber früher oder später werden sie doch...“

„Psst. Sie kommen!“

„Was sollen wir tun?“

„Nichts. Warte.“

„Aber wir können nicht für immer im Verborgenen bleiben. Wir müssen uns offenbaren!“

„Nicht jetzt...!“

„Wann dann?“

„Nicht jetzt! Bleib im Schatten!“

So warteten Thorald und Will im Heuwagen, während sich die Tore von Nottingham Castle öffneten und Wachpatrouillen förmlich hinausströmten. Die wütende Stimme des Prinzen würde Thorald bestimmt noch lange im Ohr bleiben: „Findet mir die Geächteten! Und bringt mir diesen Thorald von Andor!“

Thorald von Andor saß indes im Heuwagen direkt unter der Nase des Prinzen und fühlte sich trotz seiner seltsamen Umgebung zum ersten Mal seit langem wieder wie ein König. Er hatte Ken getrotzt und war seinen eigenen Wünschen gefolgt. Er hatte Gutes getan und einen unschuldigen Mann aus den Fängen eines Tyrannen befreit. Und das fühlte sich verflixt gut an.\bigskip







Kurz vor Sonnenaufgang wagten Thorald und Will, den Heuwagen zu verlassen und sich in Richtung des Sherwood zu begeben. Dort wartete auch schon Ken im Schatten und beglückwünschte die beiden zur gelungenen Flucht. Will führte die beiden Andori auf eine Baumbrücke und die drei Fliehenden kletterten sorgsam durch die Baumkronen des Sherwood, während unter ihnen planlose Wachen die Waldpfade entlangschlichen.

Lange dauerte es nicht, bis sie auf Robin und Marian stießen. Die beiden unterhielten sich gerade über die Chance, dass im Laufe des nächsten Tages eine Kutsche durch den Sherwood Forest reisen würde.

„Länger als zwei Tage können wir nicht warten, sonst ist es definitiv um Will geschehen“, sprach Marian gerade, „Wir sollten uns auf die Suche nach einem Seil machen und die Kutschen Kutschen sein lassen.“

„Schon seit Tagen hat keine Kutsche mehr den Sherwood passiert, jetzt muss praktisch eine auftauchen. Wir kennen ihre Strecken und wir können sie abpassen“, entgegnete Robin, „Mit John an unserer Seite sollten wir den Kutscher problemlos überwältigen können. Das Dorf liegt weit weg und ein Seil zu finden könnte zu lange dauern. Vielleicht sind die Kutschen unsere letzte Hoffnung.“

Marian schüttelte ihren Kopf: „Du und John könnt die Kutschen abpassen. Ich schleiche mich lieber ins Dorf. Bei der aktuellen Flussströmung bin ich mit dem Boot so schnell bei Megan und Alex, da wird die Sonne noch nicht einmal aufgegangen sein. Nur um Gisi müssen wir uns Sorgen machen...“

„Verkleide dich. Oder nimm die Schilde mit“, schlug Robin vor und hantierte an einem Sack auf seinem Rücken herum, „Dieser himmelblaue Dreiecksschild scheint mir besonders wirkungsvoll. Ich kann nicht bestimmen, warum dies so ist, aber in seiner Präsenz fühle ich mich sicherer. Und nicht nur mich. Nenn mich abergläubisch, aber seitdem wir diese Schilde herumtragen, ist die Hoffnung im Land unnatürlich schnell gestiegen.“

Robin kniete sich nieder und öffnete den Sack. Unüberraschenderweise enthüllte er den Sternenschild und den Bruderschild, die beiden mächtigen Schilde, die Little John Thorald kurz nach seiner Ankunft in England abgeknüpft hatte.

„Wie gerufen“, lächelte Ken, und sprang zu den beiden hinüber. Thorald und Will taten es ihm gleich.

Thorald erkannte, wie die Gesichter von Robin und Marian rasch zahlreiche Emotionen durchliefen. Schrecken, als die drei Gestalten auf ihre Baumbrücke sprangen. Entspannung, als sie Ken und Thorald erkannten, kurz gefolgt von Überraschung, als der geschundene Will dazukam. Erleichterung, als sie schlussfolgerten, dass die beiden Andori Will auf eigene Faust befreit hatten. Dann wieder Schrecken und Wut. Moment mal, Schrecken und Wut?

Jetzt war es an Thorald, überrascht zu sein. Während sich Robin und Marian langsam und vorsichtig aufrichteten, blickte sich Thorald ebenso vorsichtig um und erkannte die Quelle ihres Zorns: Ken Dorr hatte sein Schwert gezückt und hielt es dem verletzten Will Scarlet an die Kehle.

„Beim Barte des Urtrolls, was tust du da?!“, zischte Thorald Ken zu. Ken ignorierte ihn und hielt Will weiterhin fest im Griff, während er das Wort an Robin und Marian richtete: „Geächtete! Wir haben euren Will für euch aus dem Castle befreit und werden euch nun für immer verlassen. Wir hätten aber vorher gerne unsere Schilde zurück. Wenn ihr also so freundlich wärt...“

„Ken, das ist doch Wahnsinn!“, sprach Thorald. Als Ken ihn weiterhin ignorierte, wandte sich Thorald stattdessen an Robin und Marian: „Lasst nicht zu, dass diese Situation hier eskaliert. Wir... wir wollen keinen Ärger. Ken meint das nicht so, er ist nur... nur ein sehr vorsichtiger Mensch. Bitte... gebt uns einfach die beiden Schilde und wir sind auf und davon und Will ist sicher.“

„Und unser Pferd wollen wir auch zurück!“, schnarrte Ken.

„Vergiss doch das verflixte Pferd!“, fauchte Thorald zurück. Flehend blickte er die beiden Geächteten an. Wie er befürchtet hatte, hatte Robin bereits seinen Bogen und Marian ihre Axt fest gepackt und starrten Ken hasserfüllt an. Ken deeskalierte die Angelegenheit nicht im Geringsten, als er sein Schwert Will noch stärker an den Hals drückte und „Die Schilde, aber schnell!“ rief. Will wimmerte.

Ein Glück, dass wenigstens Robin und Marian noch Verstand in der Birne hatten. Marian ließ ihre Axt fallen und nickte: „Natürlich, nehmt die beiden Schilde. Sie bedeuten uns nichts.“

„Na los, Thorald, greif sie dir!”, befahl Ken.

Robin blieb angespannt, griff aber nicht ein, als Thorald, die Hände beschwichtigend gehoben, nach vorne lief und die beiden mächtigen Schilde einpackte. Sein Herz schlug rasend schnell und er achtete darauf, keine hektischen Bewegungen zu machen. Marian tat es ihm gleich. Robin schien sich ganz und gar nicht zu bewegen, sondern hatte seine Augen starr auf Ken gerichtet.

Schließlich war die Schildübergabe abgeschlossen. Ken befahl den beiden Geächteten, ihre Waffen auf den Waldboden zu werfen und dann weit zurückzutreten. Robin und Marian handelten wie angewiesen. Da ließ Ken Will frei, stieß ihn nach vorne, sprang selbst auf den Waldboden und sauste davon. Thorald hatte sich im Castle etwas von seinen Anstrengungen erholen können, warf den drei Geächteten einen letzten entschuldigenden Blick zu und sprang dann selbst auf den Waldboden.

Keine Wachen in Sicht. Auch kein Guy von Gisbourne. Die mächtigen Schilde schleppte er auf dem Rücken. Der Weg nach Andor kannte er. Sein Pferd hatte er verloren und ein paar neue Feinde hatte er sich gemacht, aber abgesehen davon war nun alles bereit für die Rückkehr nach Andor.

So nahm Thorald seine Beine in die Hand und setzte Ken nach.\bigskip







Als Thorald in die Nähe des Wasserfalls trat, hinter dem ein steinernes Gesicht den Höhlengang nach Andor verschloss, erklang ein vertrautes leises Summen aus seiner Tasche. Er wühlte noch während des Rennens darin und sah mit Freude, wie der kleine grüne Runenstein immer heller zu leuchten begann, je näher sie dem Höhleneingang kamen.

Als Thorald dem Wasserfall den Runenstein entgegenhielt, knirschte es dahinter laut. Das dort liegende Gesicht hatte offenbar gerade seinen steinernen Mund geöffnet. Die beiden Andori traten raschen durchs kühle Nass in den Gang, der sie hoffentlich nach Andor zurückbringen konnte.

\az{Jahr 74}

Die Rückkehr verlief ziemlich ereignislos. Pferdelos trotteten der König und sein engster Ratgeber durch den schummrig erleuchteten Gang und ließen die Ereignisse des vergangenen Tages Revue passieren. Zumindest tat Thorald das, und vermutete, dass es Ken ähnlich ging. Immerhin kam es nicht jeden Tag vor, dass man in eine fremde Welt gelangte und dort ein buchstäbliches Beispiel eines tyrannischen Königs unter die Nase gehalten kriegte. Vielleicht hatte das Schicksal doch noch Pläne für Thorald und versuchte ihm so klarzumachen, dass es doch nicht das übelste wäre, die mächtigen Schilde für sich zu behalten? Immerhin schadete er seinem Volk nicht absichtlich. Wer konnte schon wissen, was Krams Nachfolger mit den mächtigen Schilden anstellen könnte, sollten etwa die „Wahren Schildzwerge“ Kontrolle über Cavern erringen. Andererseits konnte es auch nicht so weitergehen wie bislang, oder Ken würde Thorald mit Worten und den Schilden immer weiter manipulieren.

Wenn er also die mächtigen Schilde nicht aus seinem Leben entfernen würde, dann...

So traf Thorald schweren Herzens eine Entscheidung. Er brach das Schweigen und seine Stimme hallte durch den breiten Stollen: „Die mächtigen Schilde sind genau das, mächtig. Ich mag nicht geeignet sein zum König, aber immerhin bin ich nicht bösartig. Das wissen wir. Von den Zwergen wissen wir das nicht. Nicht auszudenken, was der Sternenschild in den falschen Händen anrichten könnte. Es ist meine Pflicht, sie bei mir zu behalten.“

Ken blickte ihn für einen Augenblick überrascht an und schien etwas entgegnen zu wollen, zuckte dann aber mit den Schultern und nickte.

Thorald fragte sich, was in seinem Kopf vorging. Aber er wusste besser, als nachzufragen – Ken vertrug solche Fragen nicht sonderlich gut. Thorald haderte selbst noch damit, die richtigen Worte zu finden, und so kam es, dass die beiden Andori bereits erschöpft durch die Höhle bei der Quelle des Likko gestrauchelt und aus dem großen Wasserfall auf der Andor-Seite des unterirdischen Gangs getreten waren, als Thorald auffiel, dass Ken ziemlich sicher noch auf eine Entschuldigung dafür wartete, dass Thorald mit den mächtigen Schilden vor ihm davongerannt war.

„Entschuldige, Ken“, murmelte er leise.

Ken lachte nur auf: „Du ziehst alleine aus der Rietburg los, reitest mir davon, bringst uns alle in diese Bredouille und denkst, eine einfache Entschuldigung würde das alles in Ordnung bringen?“

Thorald schüttelte seinen Kopf: „Nein, Ken. Ich... ich habe so lange versucht, mich in deinen Augen als würdig zu erweisen und du behandelst mich immer mehr wie der letzte Dreck. Nicht nur mich, sondern auch alle anderen. Was hast du dir nur dabei gedacht, die Geächteten zu bedrohen? Wir haben Glück, dass Robin dir nicht gleich zwei Pfeile in den Rücken gejagt hat. Und ich... ich bin deine Manipulationen und Ausbrüche einfach satt. Du bist ein gescheiter Ratgeber, aber... aber du bist kein guter Partner. Ich glaube, dass es... dass es mir guttun würde, mich zumindest eine Zeit lang von dir fernzuhalten. Und ich entschuldige mich dafür, dass ich das wohl tun werde.“

Ken packte Thorald an den Schultern und sah ihn mit großen Augen an: „Gute Güte, Thorald, woher kommen denn all diese Ideen auf einmal? Ich bin immer noch den treuster Ratgeber und deine beste Hilfe! Du brauchst mich!“

„Ich kann auch andere Ratgeber finden, die mich besser behandeln. Elga war das einmal, bis du ins Bild kamst. Und ich brauche dich erst recht nicht, um mich zu kontrollieren und meine Fehler im Schach zu halten. Ich habe Will befreit und das Einzige, was dabei schiefgelaufen ist, war dein brutales Eingreifen.“

Ein Anflug von Furcht flackerte in Kens Augen auf, wurde dann aber gleich wieder von Wut überdeckt: „Sei kein tumber Trottel von einem Troll, Thorald! Meinst du, eine einzige gute Tat würde all deine Schlampereien gutmachen und zukünftige Fehler abwehren? Ganz abgesehen davon konntest du Will nur mit meinem Casamatuc befreien, und nur weil ich den Karren mit dem Heu vor die Burg geschoben habe! Du brauchst mich! Deinen treuen Ratgeber, deinen Diener, deinen ergebenen Ken Dorr!“

Kens letzte Worte hatten schon beinahe trotzig geklungen.

Thorald stockte, fuhr dann aber fort: „Nein, Ken, es ist gerade andersherum: Du brauchst mich, Ken. Du bist von meiner Macht abhängig, und du hast zu lange versucht, mich zu kontrollieren, und mein Wohlbefinden aktiv missachtet. Nicht weiter. Ich ziehe jetzt zurück zur Rietburg, lagere die mächtigen Schilde dort und bleibe gerade so lange im Amt, bis die Helden aus dem Norden zurückkehren und ich ihnen das Amt übertragen kann. Ich habe viele Fehler gemacht und vielen das Leben gekostet. Aber diese Entscheidung ist kein Fehler. Manche... manche Leute sind einfach nicht für das Königsamt gemacht. Ich mag einer davon sein. Du bist das aber ebenfalls.“

Thorald blieb einen kurzen Moment ergriffen stehen und genoss das Licht der andorischen Sonne auf seinem Gesicht. Die Luft roch einfach anders hier. In einiger Ferne sah er Elgas und Junas Bauernhof am Fuße des Grauen Gebirges stehen. Wahrscheinlich lag die Rietgraskrone immer noch irgendwo dort herum. Thorald ließ seinen Blick über die nebelumwobene Narne hinweg zur Rietburg schweifen. Lieber direkt zur Burg zurückkehren und die Schilde in Sicherheit bringen. Der König von Andor würde von einem Abenteuer zurückkehren und zufrieden sein. Die Rietgraskrone konnte er immer noch später abholen. Bei Elga und Juna war sie wahrscheinlich sogar sicherer als in der Burg.

„Ich habe Freunde in der Burg, viele Freunde“, zischte Ken und unterbrach so Thoralds Gedanken, „Ich habe in meiner Zeit als Soldat des Königshauses Seite an Seite mit vielen Männern der Rietgarde gekämpft, während du über ihren Köpfen Wein gesoffen hast! Viele waren erzürnt über meine Verbannung durch das Hause Brandurs, das du repräsentierst!“

Kens Stimme hatte einen finsteren Unterton angenommen.

„Was meinst du, Thorald, auf wessen Seite sich die Garde stellen würde, sollte ich mich öffentlich gegen dich stellen? Du willst es nicht darauf ankommen lassen!“

Thorald verwarf die Drohung: „Du willst das doch ebenso wenig. Die Schildzwerge würden sich nur liebend gerne unter einem Vorwand auf unsere mächtigen Schilde stürzen. Eine Revolution käme ihnen nur allzu gerne recht. Und die Helden... die Helden würden dich nicht auf dem Thron stehen lassen. Das ist immer noch ihr Land.“

„Du begehst einen großen Fehler, Thorald!“

„Ich ziehe von hinnen, Ken. Es ist vorbei.“

Thorald wandte sich energisch von Ken ab und stapfte in Richtung der Bogenbrücke, die beiden mächtigen Schilde auf seinem Rücken tragend. Er fühlte sich... erleichtert, irgendwie. Als wäre ein großes Gewicht von seinen Schultern genommen worden. Der Sternenschild auf seinem Rücken summte beruhigend. Alles war gut.








\newpage
\section{Epilog}

„Du brauchst mich!“, rief Ken Thorald ein letztes Mal hinterher. Thoralds Gestalt wurde immer kleiner, während der Prinz davonschritt und seinen Ken nicht weiter beachtete. Ken stampfte mit seinem Fuß auf, schritt ihm aber nicht hinterher. Er durfte nun nicht planlos handeln. Sein Kopf ratterte, während sein Blick noch lange auf seinem König verblieb.

Schritte.

„Da seid ihr ja endlich!“, knurrte Ken seine beiden Handlanger aus der Rietgarde an, welche keuchend angerannt kamen, „Ihr beiden habt ja was verpasst.“

„Wo seid Ihr denn gewesen, Herr?“, frage der eine.

„War das der König, den ich soeben abstürmen gesehen habe?“, fragte der andere.

Ken Kiefer zitterte, als er die Worte ausstieß: „Ich... ich werde ihm nachgehen. Ich werde ihn wieder umstimmen. Ich kann das. Er... er muss auf mich hören, alleine ist er der Krone nie und nimmer gewachsen... er begeht gerade den größten Fehler seines Lebens... und wenn er das bald nicht eingesehen hat, dann... dann..., wenn nicht. Aber er wird es tun.“

Ken knurrte weiter vor sich hin, und seine beiden Handlanger tauschten nervöse Blicke aus, blieben aber stumm. Als Ken das Wort wieder ergriff, triefte seine Stimme nur so von Zorn und Bosheit.

„Thorald darf auf keinen Fall die Rietburg erreichen und die Stadtwache gegen mich aufbringen. Wir müssen ihn irgendwie ablenken, seine Wut auf etwas anderes weisen. Elga und Juna, diese Bäuerinnen am Fuße des Gebirges, deren Gesellschaft Thorald der meinen so sehr vorzuziehen scheint – sie haben die Rietgraskrone noch. Holt sie euch und bringt die Krone zur Burg. Sollten den Bäuerinnen unterwegs ein... Unfall zustoßen, dann wäre das natürlich ganz tragisch.“

Die leere Blick seiner Handlanger klärte sich. Mit solchen Worten wussten sie schon besser umzugehen.

„Nun stellt euch vor, dass jemand mit einer blutigen Krone ins Rietland zurückkehrte... da könnte der Prinz ja auf die Idee kommen, dass die räudigen Wölfe des Südens Elga und Juna erwischt hätten. Er, der stets nach großen Taten dürstet, würde seinen Hass auf die Wölfe lenken, wieder auf meinen Rat hören und diesem folgend zur Jagd blasen lassen. Diese sollte ihn von mir ablenken und wieder etwas in die Realität zurückholen. Danach würde er hoffentlich wieder auf meinen Rat hören. Und falls er sich weiter gegen mich wehrte... Nun, wenn ein König nicht weiß, was gut für sein Land ist, wenn er eine Gefahr für sich selbst und seine Gemeinschaft anstellt, dann muss man sich seiner entledigen. Ich hatte gehofft, dass es nie dazu kommen würde, aber ich bin ein Realist. Ich werde euch entsprechende Anweisungen zukommen lassen. Falls niemand in der Burg von meinem Zerwürfnis mit Thorald erführe und sein Dahinscheiden natürlich wirkte, so würde keiner in der Burg es wagen, meine Position als Statthalter von Andor in Frage zu stellen. Aber zunächst einmal gilt es, den König auf eine Wolfsjagd zu hetzen. Los, los, ab jetzt. Auf zu den Bäuerinnen!“

Kens Handlanger sahen einander ein letztes Mal fragend an und zogen dann hastig von hinnen. Ken lehnte sich zurück und sog tief Atem ein. Ja, Andor roch einfach anders als dieser seltsame Wald, in dem er und Thorald gelandet waren. Dieser verfluchte Wald in diesem verfluchten fremden Land! Er würde noch einmal zur Höhle reisen müssen und den Zugang zerstören. Einstürzen lassen Zur Gänze vernichten. Nicht auszudenken, was geschähe, wenn die Armee dieses anderen Königs in Andor einfallen könnte.

Ken schüttelte sich und brach auf. Auf, zurück nach Andor. Auf, zurück zur Rietburg. Er musste seine Karten jetzt gut ausspielen, um am Ende noch obenauf zu enden. Aber ja, alles konnte noch gut kommen.\bigskip







Eine kurze Zeit lang lag die Wiese neben der Quelle des Likko still da. Aber verlassen war sie nicht. Denn sobald Ken außer Hör- und Sichtweite getrabt war, löste sich eine grün gewandete Gestalt aus dem Schatten und sah sich vorsichtig um, den Bogen stets im Anschlag.

Robin von Loksley hatte nicht alles verstanden, was dieser Kennard von Dorr mit seinen Handlangern besprochen hatte. Aber er hatte genug gehört, um zu wissen, dass Bauern in Gefahr waren. Und wie könnte er sich guten Herzens einen Helden und Beschützer der Armen nennen, wenn er nicht Menschen in Not beistehen würde, egal, aus welcher Welt sie stammten? In einer war Robin ja schon ein Geächteter. Zeit, es auch in einer zweiten zu werden. Vielleicht würde dies ja endlich seinen Albträumen ein Ende bereiten.

So wandte sich Robin um und winkte zur Quelle des Flusses zurück. Prompt schlichen seine Gefährten aus dem Schatten hervor und schlossen sich Robin an. Gemeinsam folgte die Gruppe im Geheimen Kennards Handlangern in Richtung Bauernhof.

Robin zählte fünf Pfeile in seinem Köcher. Er lächelte.

Mehr würde er auch nicht brauchen.
